\documentclass[10pt,a4paper]{amsart}
\usepackage{amsmath,amssymb,amsthm,graphicx}
\usepackage{hyperref}
\usepackage{amssymb}
\usepackage{xcolor}
\numberwithin{equation}{section}

\newtheorem{Theorem}{Theorem}[section]
\newtheorem*{Theorem*}{Theorem}
\newtheorem{Corollary}[Theorem]{Corollary}
\newtheorem{Lemma}[Theorem]{Lemma}
\newtheorem{Proposition}[Theorem]{Proposition}
 { \theoremstyle{definition}
\newtheorem{Definition}[Theorem]{Definition}
\newtheorem{Note}[Theorem]{Note}
\newtheorem{Example}[Theorem]{Example}
\newtheorem{Remark}[Theorem]{Remark}
\newtheorem{Notation}[Theorem]{Notation}}

\usepackage[all]{xy}
\newcommand{\rond}[1]{*++[o][F-]{#1}}

\newcommand{\eas}{\mathbf{EAS}}
\newcommand{\leas}{\ell\mathbf{EAS}}
\allowdisplaybreaks
\title{Structure classification of finite non-degenerate commutative extended diassociative semigroups}
\author{Loïc Foissy}
\date{doi:10.3842/SIGMA.2025.092; CC BY 4.0}
\begin{document}
\maketitle
\noindent\textit{Allowed standard prerequisites (editorial).}\begin{itemize}
\item Research prerequisites named in the original article. Full original bibliography retained; all source-local definitions and supporting arguments retained.
\end{itemize}\medskip
\section{Introduction}

It seems that the notion of \emph{family} parameterization of a given type of algebraic structure appeared firstly appears in the context of quantum field theory:
in \cite{Ebrahimi-Fard2007}, the authors introduced Rota--Baxter family algebras. This terminology is due to Li Guo \cite{Guo2009}. In the same spirit, family pre-Lie, dendriform or tridendriform algebras, among others have been introduced~\cite{Foissy49,Zhang2019,Zhang2020,Zhang2020-3,Manchon2020}. In all cases, the idea is to replace the operations defining the structure by a~bunch of operations indexed by a~semigroup $\Omega$;
the relations between the axioms are deformed using this structure on $\Omega$.
For example, if $(\Omega,\star)$ is a semigroup, an $(\Omega,\star)$-family associative algebra $A$ has a family $(*_\alpha)_{\alpha \in \Omega}$ of products, with the relations
\[x*_\alpha (y *_\beta z)=( x*_\alpha y)*_{\alpha \star \beta}z,\]
satisfied for any $\alpha,\beta,\gamma \in \Omega$ and any $x$, $y$, $z$ in $A$. In the same spirit, the notion of \emph{matching} parameterization can be used: for pre-Lie, it appears for example in the work of Bruned, Hairer and Zambotti
on regularity structures to solve stochastics PDEs \cite{Bruned2019,Bruned2023,Bruned2023-1,Foissy47}.
Matching Rota--Baxter algebras, associative, dendriform, pre-Lie algebras are introduced in \cite{Foissy47,Zhang2020}, see also \cite{Foissy49} for a two-parameter versions for pre-Lie algebras.
For example, a matching associative algebra has a family $(*_\alpha)_{\alpha \in \Omega}$ of products indexed by a set $\Omega$, with the relations
\[x*_\alpha (y *_\beta z)=(x *_\alpha y)*_\beta z.\]
Note that no specific structure is required on $\Omega$ in this case.
Attempts to unify these parameterizations have been done in \cite{Foissy56,Foissy48,Foissy50}. For example, for associative, following \cite{Foissy56},
given a set~$\Omega$ with two binary operations $\rightarrow$ and $\triangleright$, an $(\Omega,\rightarrow,\triangleright)$-associative algebra has a family $(*_\alpha)_{\alpha \in \Omega}$ of products, with the relations
\[x*_\alpha (y *_\beta z)=x *_{\alpha\triangleright \beta}(y*_{\alpha \rightarrow \beta}z).\]
Of course, usually these structures do not have any convenient property, and some conditions are imposed: roughly speaking, one imposes that the underlying combinatorics of the initial object is conserved, modulo a parameterization by $\Omega$.
This gives some constraints on $\Omega$. For associative algebras, $(\Omega,\rightarrow,\triangleright)$ has to be an extended associative semigroup (briefly, EAS):
\begin{gather*}
\forall \alpha,\beta,\gamma \in \Omega,\qquad \alpha \rightarrow (\beta \rightarrow \gamma)=(\alpha \rightarrow \beta) \rightarrow \gamma,\\
(\alpha \triangleright (\beta \rightarrow \gamma))\rightarrow (\beta \triangleright \gamma)=(\alpha \rightarrow \beta)\triangleright \gamma,\qquad
(\alpha \triangleright (\beta \rightarrow \gamma))\triangleright (\beta \triangleright \gamma)=\alpha \triangleright \beta.
\end{gather*}
In particular, $(\Omega,\rightarrow)$ is an associative semigroup. Here are some examples:
\begin{itemize}\itemsep=0pt
\item If $\Omega$ is a set, putting
$
 \forall \alpha,\beta \in \Omega$, $ \alpha \rightarrow\beta =\beta$, $\alpha \triangleright \beta =\alpha$,
we obtain an EAS, denoted by $\eas(\Omega)$. This EAS gives back matching associative algebras.
\item If $(\Omega,\rightarrow)$ is an associative semigroup, it is an EAS with
$\forall \alpha,\beta \in \Omega$, $\alpha \triangleright \beta =\alpha$.
This EAS is denoted by $\eas(\Omega,\rightarrow)$. It gives back $(\Omega,\rightarrow)$-family associative algebras.
\item If $(\Omega,\star)$ is a group, it is an EAS, with
$
 \forall \alpha,\beta \in \Omega$, $\alpha \rightarrow\beta =\beta$, $\alpha \triangleright \beta =\alpha\star \beta^{\star -1}$,
It is denoted by $\eas'(\Omega,\star)$. \end{itemize}
The two first examples explain why $\Omega$-matching and $(\Omega,\star)$-family associative algebras are very similar, in particular why the free objects are isomorphic as vector spaces:
this is fact works for the more general settings of $\Omega$-associative algebras over an EAS.
The same can be done with pre-Lie algebras, leading to the notion of commutative extended diassociative semigroup (briefly, CEDS). A CEDS is an EAS satisfying the complementary axioms
\begin{gather*}
\forall \alpha,\beta,\gamma \in \Omega,\qquad(\alpha \rightarrow \beta) \rightarrow \gamma =(\beta \rightarrow \alpha) \rightarrow \gamma,\qquad
\alpha \triangleright (\beta \rightarrow \gamma) =\alpha \triangleright \gamma,
\end{gather*}
The Koszul duality of quadratic operads applied to $\Omega$-pre-Lie algebras leads to the notion of dual CEDS, which are EAS with the complementary axioms
\begin{gather*}
 \forall \alpha,\beta,\gamma \in \Omega,\qquad (\alpha \triangleright \beta) \rightarrow \gamma=\alpha \rightarrow \gamma,\qquad
(\alpha \triangleright \beta) \triangleright \gamma =(\alpha \triangleright \gamma)\triangleright \beta.
\end{gather*}
For example, for any set $\Omega$, $\eas(\Omega)$ is both a CEDS and a dual CEDS. For any semigroup~$(\Omega,\star)$, $\eas(\Omega,\star)$ is a dual CEDS and is a CEDS if and only if
\begin{align*}
\forall \alpha,\beta,\gamma \in \Omega,\alpha \star (\beta \star \gamma)=(\alpha \star \beta) \star \gamma=(\beta \star \alpha) \star \gamma.
\end{align*}
For any group $(\Omega,\star)$, $\eas'(\Omega,\star)$ is a CEDS, and is a dual CEDS if and only if $\star$ is commutative.

The axioms of EAS can be reformulated using the maps
\begin{gather*}
\phi\colon\ \begin{cases}
\Omega^2\longrightarrow\Omega^2,\\
(\alpha,\beta)\longrightarrow(\alpha \rightarrow \beta,\alpha \triangleright \beta),
\end{cases}\qquad
\tau\colon\ \begin{cases}
\Omega^2\longrightarrow\Omega^2,\\
(\alpha,\beta)\longrightarrow(\beta,\alpha).
\end{cases}
\end{gather*}
Then $(\Omega,\rightarrow,\triangleright)$ is an EAS if and only if
\begin{align*}
(\operatorname{Id} \times \phi)\circ (\phi\times \operatorname{Id})\circ (\operatorname{Id} \times \phi)&=(\phi \times \operatorname{Id})\circ (\operatorname{Id} \times \tau)\circ (\phi\times \operatorname{Id}).
\end{align*}
Similar formulations can be done for CEDS and dual CEDS, see Lemma \ref{lem2.6}.
This reformulation naturally leads to the notion of linear EAS: an $\ell$EAS is pair $(A,\Phi)$, where $A$ is a vector space and $\Phi\colon A\otimes A\longrightarrow A\otimes A$ is a linear map satisfying the \emph{$\ell$EAS braid equation}
\begin{align*}
(\operatorname{Id} \otimes \Phi)\circ (\Phi\otimes \operatorname{Id})\circ (\operatorname{Id} \otimes \Phi)&=(\Phi \otimes \operatorname{Id})\circ (\operatorname{Id} \otimes \tau)\circ (\Phi\otimes \operatorname{Id}).
\end{align*}

In particular, let $(\Omega,\rightarrow,\triangleright)$ be a set with two operations. We denote by $\mathbb{K}\Omega$ the vector space generated by $\Omega$ and we define $\Phi\colon\mathbb{K}\Omega\otimes \mathbb{K}\Omega\longrightarrow\mathbb{K}\Omega\otimes \mathbb{K}\Omega$ by
\begin{align*}
&\forall \alpha,\beta\in \Omega,\qquad\Phi(\alpha\otimes \beta)=(\alpha \rightarrow \beta)\otimes (\alpha \triangleright \beta).
\end{align*}
Then $(\mathbb{K}\Omega,\Phi)$ is an $\ell$EAS if and only if $(\Omega,\rightarrow,\triangleright)$ is an EAS.
Not all the $\ell$EAS can be obtained in this way, see Example~\ref{ex4.1} for two-dimensional examples. Similar presentations of $\ell$CEDS and $\ell$CEDS can be established, see Definition \ref{defi4.1}.

The aim of this paper is a study of EAS, CEDS, and their linear versions. In the second section, after recalling the main definitions, we give a classification of EAS of cardinality 2,
which gives 13 non-isomorphic examples, 11 being CEDS, 7 being dual CEDS, 3 being non-degenerate (that is to say, with an invertible map $\phi$).
The third section is devoted to the study of non-degenerate finite CEDS. We prove that the three examples defined earlier are in fact fundamental bricks:
Theorem~\ref{theo3.14} states that any finite non-degenerate CEDS can be decomposed as the direct product of a semi-direct product $\eas(\Omega_1,*) \rtimes \eas'(\Omega_2,\star)$ with an~$\eas(\Omega_3)$,
where $(\Omega_1,*)$ is an abelian group, $(\Omega_2,\star)$ is a group and $\Omega_3$ is a nonempty set. The fourth section is devoted to linear versions of EAS.
We give firstly a family of 18 examples of~$\ell$EAS in dimension 2, then study the duality of $\ell$EAS (see Proposition \ref{prop4.2}), and left units, left counits and eigenvectors (see Definition \ref{defi4.3}).
If $(A,\Phi)$ is an $\ell$EAS, an element $a\in A$ is a~left unit if for any $b\in A$, $\Phi(a\otimes b)=b\otimes a$. An element $f \in A^*$ is a~left counit if for any $a,b\in A$, $(f\otimes \operatorname{Id})\circ \Phi(a\otimes b)=f(b)a$.
In particular, we characterize left units and counits and eigenvectors for linearization of non-degenerate CEDS in Proposition \ref{prop4.5}.
In the last section, we introduce two functors taking their values in the category of $\ell$EAS.
The first one (see Proposition \ref{prop5.1}) is defined on the category of bialgebras (not necessarily unitary nor counitary) and generalizes the construction of $\eas(\Omega,\rightarrow)$.
The second one (see Proposition \ref{prop5.7}) is defined on the category of Hopf algebras and generalizes the construction of $\eas'(\Omega,*)$.
In the case of an $\ell$EAS coming from a Hopf algebra, this is closely related to the notion of right integral (see Proposition \ref{prop5.11}).
We prove in Theorem~\ref{theo5.12} that we can associate to any convenient pair $(a,f)$ of a unit and a~counit a~bialgebra structure on $A$, recovering in this way $\ell$EAS coming from a bialgebra.
This is finally applied to $\ell$EAS defined from Hopf algebras of groups.



\begin{Notation}
$\mathbb{K}$ is a commutative field. All the vector spaces in this text are taken over $\mathbb{K}$.
\end{Notation}


\section{Extended (di)associative semigroups}

\subsection{Commutative extended diassociative semigroups}

Let us first recall this definition of \cite{Foissy49}, where it is related to a parameterization of the operad of dendriform algebras.

\begin{Definition}\qquad
\begin{enumerate}\itemsep=0pt
\item[(1)] A \emph{diassociative semigroup} is a family $(\Omega,\leftarrow,\rightarrow)$, where $\Omega$ is a nonempty set and $\leftarrow$, $\rightarrow\colon\Omega\times \Omega \longrightarrow \Omega$ are maps such that, for any $\alpha,\beta,\gamma\in \Omega$,
\begin{gather}%\label{eq2}\label{eq3}
\label{eq1} (\alpha \leftarrow \beta)\leftarrow \gamma=\alpha\leftarrow (\beta \leftarrow \gamma)=\alpha \leftarrow(\beta \rightarrow \gamma),\\
 (\alpha \rightarrow \beta)\leftarrow \gamma=\alpha \rightarrow (\beta \leftarrow \gamma),\\
 (\alpha \rightarrow \beta)\rightarrow \gamma=(\alpha \leftarrow \beta) \rightarrow \gamma=\alpha \rightarrow (\beta \rightarrow \gamma).
\end{gather}
An \emph{extended diassociative semigroup} (briefly, EDS) is a family $(\Omega,\leftarrow,\rightarrow,\triangleleft,\triangleright)$, where $\Omega$ is a nonempty set
and $\leftarrow,\rightarrow,\triangleleft,\triangleright\colon\Omega\times \Omega \longrightarrow \Omega$ are maps such that
\begin{enumerate}\itemsep=0pt
\item $(\Omega,\leftarrow,\rightarrow)$ is a diassociative semigroup.
\item For any $\alpha,\beta,\gamma \in \Omega$,
\begin{gather}%\label{eq4}\label{eq5}\label{eq6}\label{eq7}\label{eq8}\label{eq10}\label{eq11}
\alpha\triangleright (\beta \leftarrow \gamma)=\alpha \triangleright \beta,\\
(\alpha \rightarrow \beta)\triangleleft \gamma=\beta \triangleleft \gamma, \\
(\alpha \triangleleft \beta)\leftarrow ((\alpha \leftarrow \beta) \triangleleft \gamma)=\alpha \triangleleft (\beta \leftarrow \gamma),\\
(\alpha \triangleleft \beta)\triangleleft ((\alpha \leftarrow \beta) \triangleleft \gamma)=\beta \triangleleft \gamma,
\\
(\alpha \triangleleft \beta) \rightarrow ((\alpha \leftarrow \beta) \triangleleft \gamma)=\alpha \triangleleft (\beta \rightarrow \gamma),\\
 (\alpha \triangleleft \beta) \triangleright ((\alpha \leftarrow \beta) \triangleleft \gamma)=\beta \triangleright \gamma,\\
(\alpha \triangleright (\beta \rightarrow \gamma))\leftarrow (\beta \triangleright \gamma)=(\alpha \leftarrow \beta) \triangleright \gamma,\\
(\alpha \triangleright (\beta \rightarrow \gamma))\triangleleft (\beta \triangleright \gamma)=\alpha \triangleleft \beta,\\
\label{eq12} (\alpha \triangleright (\beta \rightarrow \gamma))\rightarrow (\beta \triangleright \gamma)=(\alpha \rightarrow \beta)\triangleright \gamma,\\
\label{eq13} (\alpha \triangleright (\beta \rightarrow \gamma))\triangleright (\beta \triangleright \gamma)=\alpha \triangleright \beta.
\end{gather}
\end{enumerate}
\end{enumerate}
\end{Definition}

An EDS $(\Omega,\leftarrow,\rightarrow,\triangleleft,\triangleright)$ is commutative if for any $\alpha,\beta \in \Omega$,
\begin{align}
\label{eq14}\alpha \leftarrow \beta=\beta \rightarrow \alpha,\qquad\alpha \triangleleft \beta=\beta \triangleright \alpha.
\end{align}

Let us reformulate the definition of commutative EDS.

\begin{Proposition}
A commutative EDS $($briefly, CEDS$)$ is a triple $(\Omega,\rightarrow,\triangleright)$, where $\Omega$ is a~non\-empty set and $\rightarrow, \triangleright\colon\Omega^2\longrightarrow \Omega$ are maps such that, for any $\alpha,\beta,\gamma \in \Omega$,
\begin{gather}
\label{eq15} \alpha \rightarrow (\beta \rightarrow \gamma)=(\alpha \rightarrow \beta) \rightarrow \gamma=(\beta \rightarrow \alpha) \rightarrow \gamma,\\
\label{eq16} \alpha \triangleright (\beta \rightarrow \gamma)=\alpha \triangleright \gamma,\\
\label{eq17} (\alpha \triangleright \gamma)\rightarrow (\beta \triangleright \gamma)=(\alpha \rightarrow \beta)\triangleright \gamma,\\
\label{eq18} (\alpha \triangleright \gamma)\triangleright (\beta \triangleright \gamma)=\alpha \triangleright \beta.
\end{gather}\end{Proposition}

\begin{proof}
Replacing $\leftarrow$ and $\triangleleft$ in \eqref{eq1}--\eqref{eq13} with the help of \eqref{eq14}, we find \eqref{eq15}--\eqref{eq18}.
\end{proof}

\begin{Definition}[{\cite{Foissy56}}]
An \emph{extended associative semigroup} (briefly, EAS) is a triple $(\Omega,\rightarrow,\triangleright)$, where~$\Omega$ is a nonempty set and $\rightarrow,\triangleright\colon\Omega^2\longrightarrow \Omega$ are maps such that, for any $\alpha,\beta,\gamma \in \Omega$,
\begin{gather}
\label{eq19} \alpha \rightarrow (\beta \rightarrow \gamma) =(\alpha \rightarrow \beta) \rightarrow \gamma,\\
\tag{\ref{eq12}}(\alpha \triangleright (\beta \rightarrow \gamma))\rightarrow (\beta \triangleright \gamma) =(\alpha \rightarrow \beta)\triangleright \gamma,\\
\tag{\ref{eq13}}(\alpha \triangleright (\beta \rightarrow \gamma))\triangleright (\beta \triangleright \gamma) =\alpha \triangleright \beta.
\end{gather}\end{Definition}

\begin{Remark}
Let $(\Omega,\rightarrow,\leftarrow,\triangleright,\triangleleft)$ be an EDS. Then $(\Omega,\rightarrow,\triangleright)$ is an EAS, called the \emph{right part} of the EDS $(\Omega,\rightarrow,\leftarrow,\triangleright,\triangleleft)$.
We obtain a commutative triangle of functors
\[\xymatrix{\mathbf{CEDS}\ar@{^(->}[r]\ar@{^(->}[d]&\mathbf{EAS}.\\
\mathbf{EDS}\ar[ru]_{\text{\scriptsize{right part}}}&}\]
We shall see that not all the EAS are right parts of an EDS (see case \textbf{C6} in the classification of EAS of cardinality 2 in the next paragraph).
\end{Remark}

\begin{Example}\qquad
\begin{enumerate}\itemsep=0pt
\item[(1)] Let $\Omega$ be a set. We put
\begin{align*}
\forall (\alpha,\beta) \in \Omega^2,\qquad\begin{cases}
 \alpha \rightarrow \beta=\beta,\\
\alpha \triangleright \beta=\alpha.
\end{cases} \end{align*}
Then $(\Omega,\rightarrow,\triangleright)$ is an EAS, denoted by $\eas(\Omega)$. It is a CEDS.

\item[(2)] Let $(\Omega,\star)$ be an associative semigroup and let $\pi\colon\Omega\longrightarrow\Omega$ be an endomorphism of $(\Omega,\star)$ such that $\pi^2=\pi$. We put
$
\forall (\alpha,\beta) \in \Omega^2$, $ \alpha \triangleright \beta=\pi(\alpha)$.
It is an EAS, which we denote by~${\eas(\Omega,\star,\pi)}$. It is a CEDS if and only if for any~${\alpha,\beta,\gamma \in \Omega}$,
$(\alpha \star \beta)\star \gamma=(\beta \star \alpha)\star \gamma$.
We shall simply denote $\eas(\Omega,\star)$ instead of $\eas(\Omega,\star,\operatorname{Id}_\Omega)$. In particular, if $(\Omega,\star)$ is a~group, then $\eas(\Omega,\star)$ is a CEDS if and only if $(\Omega,\star)$ is abelian, which proves that not all EAS are CEDS.

\item[(3)] Let $\Omega$ be a set with an operation $\triangleright$ such that, for any $\alpha,\beta,\gamma\in \Omega$,
$(\alpha \triangleright \gamma)\triangleright (\beta\triangleright \gamma)=\alpha \triangleright \beta$.
We then put
$\forall (\alpha,\beta)\in \Omega^2$, $\alpha\rightarrow \beta =\beta$.
Then $(\Omega,\rightarrow,\triangleright)$ is a CEDS (so is an EAS). This holds, for example, if $(\Omega,\star)$ is an associative semigroup with the \emph{right inverse condition}
\begin{align*}
&\forall (\beta,\gamma)\in \Omega^2, \ \exists! \alpha \in \Omega,\qquad \alpha\star \beta=\gamma.
\end{align*}
This unique $\alpha$ is denoted by $\gamma \triangleright \beta$. Then, for any $\alpha,\beta,\gamma \in \Omega$,
\begin{align*}
((\alpha \triangleright \gamma)\triangleright (\beta\triangleright \gamma))\star \beta&=((\alpha \triangleright \gamma)\triangleright (\beta\triangleright \gamma))\star ((\beta \triangleright \gamma)\star \gamma)\\
&=(((\alpha \triangleright \gamma)\triangleright (\beta\triangleright \gamma))\star (\beta \triangleright \gamma))\star \gamma
=(\alpha \triangleright \gamma)\star \gamma=\alpha,
\end{align*}
so $(\alpha \triangleright \gamma)\triangleright (\beta\triangleright \gamma)=\alpha \triangleright \beta$. This EAS is denoted by $\eas'(\Omega,\star)$.
The right inverse condition holds for example if $(\Omega,\star)$ is a group, and then
$\alpha \triangleright \beta=\alpha \star \beta^{\star-1}$.
It also holds for semigroups which are not groups. For example, if $\Omega$ is a nonempty set, we give it an associative product defined by
$\forall \alpha,\beta\in \Omega$, $ \alpha \star \beta=\alpha$.
It satisfies the right inverse condition and, for any $\alpha,\beta \in \Omega$, $\alpha \triangleright \beta=\alpha$. Note that for this example, $\eas'(\Omega,\star)=\eas(\Omega)$.
\end{enumerate}\end{Example}

\begin{Definition} \label{defi2.4}
Let $(\Omega,\rightarrow,\triangleright)$ be an EAS. We shall say that it is \emph{non-degenerate} if the following map is bijective:
\[\phi\colon \begin{cases}
\Omega^2\longrightarrow\Omega^2,\\
(\alpha,\beta)\longrightarrow(\alpha\rightarrow\beta,\alpha \triangleright\beta).
\end{cases}\]
\end{Definition}

If $\Omega$ is a non-degenerate EAS, the structure implied on $\Omega$ by $\phi^{-1}$ will be studied in the next paragraph.

\begin{Example}\qquad
\begin{enumerate}\itemsep=0pt
\item[(1)] Let $\Omega$ be a set. In $\eas(\Omega)$, for any $\alpha,\beta \in \Omega$, $\phi(\alpha,\beta)=(\beta,\alpha)$, so $\eas(\Omega)$ is non-degenerate.
\item[(2)] Let $(\Omega,\star)$ be a group. Then $\eas(\Omega,\star)$ is non-degenerate. Indeed, in this case, $\phi(\alpha,\beta)=(\alpha\star\beta,\alpha)$, so $\phi$ is a bijection, of inverse given by $\phi^{-1}(\alpha,\beta)=\bigl(\beta,\beta^{\star-1}\star \alpha\bigr)$.
\item[(3)] Let $(\Omega,\star)$ be an associative semigroup with the right inverse condition. Then $\eas'(\Omega,\star)$ is non-degenerate.
Indeed, in this case, $\phi(\alpha,\beta)=(\beta,\alpha\triangleright \beta)$, so $\phi$ is a bijection, of inverse given by $\phi^{-1}(\alpha,\beta)=(\beta\star \alpha,\alpha)$.
\end{enumerate}\end{Example}


\subsection{Dual commutative extended semigroups}

\begin{Definition}
Let $(\Omega,\rightarrow,\triangleright)$ be a set with two binary operations. We shall say that it is a~\emph{dual CEDS} if, for any $\alpha,\beta,\gamma \in \Omega$,
\begin{gather}%\label{eq21}\label{eq22}
\tag{\ref{eq19}} (\alpha \rightarrow \beta)\rightarrow \gamma=\alpha \rightarrow (\beta \rightarrow \gamma),\\
\tag{\ref{eq12}}(\alpha \triangleright (\beta \rightarrow \gamma))\rightarrow (\beta \triangleright \gamma)=(\alpha \rightarrow \beta) \triangleright \gamma,\\
\tag{\ref{eq13}} (\alpha \triangleright (\beta \rightarrow \gamma))\triangleright (\beta \triangleright \gamma)=\alpha \triangleright \beta,\\
 (\alpha \triangleright \beta) \rightarrow \gamma=\alpha \rightarrow \gamma,\\
 (\alpha \triangleright \beta) \triangleright \gamma=(\alpha \triangleright \gamma)\triangleright \beta.
\end{gather}\end{Definition}

\begin{Example}\qquad
\begin{enumerate}\itemsep=0pt
 \item[(1)] If $\Omega$ is a set, then $\eas(\Omega)$ is a dual CEDS.
 \item[(2)] If $(\Omega,\star)$ is a semigroup and $\pi\colon\Omega\longrightarrow \Omega$ is a semigroup morphism such that $\pi^2=\pi$, then~${\eas(\Omega,\star,\pi)}$ is a dual CEDS if and only if
$
 \forall \alpha,\beta \in \Omega$, $\pi(\alpha)\star \beta =\alpha \star \beta$.
 In particular, $\eas(\Omega,\star)$ is a dual CEDS.
 \item[(3)] If $(\Omega,\star)$ is a semigroup with the right inverse condition, then $\eas'(\Omega,\star)$ is a dual CEDS if and only if
$\forall \alpha,\beta,\gamma \in \Omega$, $(\alpha \triangleright \beta) \triangleright \gamma =(\alpha \triangleright \gamma) \triangleright \beta$.
 This is equivalent to
$\forall \alpha,\beta,\gamma \in \Omega$, $\alpha \star \beta \star \gamma =\alpha \star \gamma \star \beta$.
 In the case where $(\Omega,\star)$ is a group, $\eas'(\Omega,\star)$ is a dual CEDS if and only if $(\Omega,\star)$ is abelian.
 \end{enumerate}\end{Example}

The following lemma, proved in \cite{Foissy48}, is a reformulation of the axioms of EAS, CEDS and dual CEDS with the help of the map $\phi$.

\begin{Lemma}\label{lem2.6}
Let $(\Omega,\rightarrow,\triangleright)$ be a set with two binary operations. We consider the maps
\begin{align*}
\phi\colon\ \begin{cases}
\Omega^2 \longrightarrow \Omega^2,\\
(\alpha,\beta) \longrightarrow (\alpha \rightarrow \beta,\alpha \triangleright \beta),
\end{cases} \qquad
\tau\colon\ \begin{cases}
\Omega^2 \longrightarrow \Omega^2,\\
(\alpha,\beta) \longrightarrow (\beta,\alpha).
\end{cases}
\end{align*}
Then
\begin{enumerate}\itemsep=0pt
\item[$(1)$] $(\Omega,\rightarrow,\triangleright)$ is an EAS if and only if
\begin{align}
\label{eq23}(\operatorname{Id} \times \phi)\circ (\phi\times \operatorname{Id})\circ (\operatorname{Id} \times \phi) =(\phi \times \operatorname{Id})\circ (\operatorname{Id} \times \tau)\circ (\phi\times \operatorname{Id}).
\end{align}
\item[$(2)$] $(\Omega,\rightarrow,\triangleright)$ is a CEDS if and only if
\begin{gather}
(\operatorname{Id} \times \phi)\circ (\phi\times \operatorname{Id})\circ (\operatorname{Id} \times \phi) =(\phi \times \operatorname{Id})\circ (\operatorname{Id} \times \tau)\circ (\phi\times \operatorname{Id}),\tag{\ref{eq23}}\\
(\operatorname{Id} \times \phi)\circ (\operatorname{Id} \times \tau) \circ (\tau \times \operatorname{Id})\circ (\phi \times \operatorname{Id})\nonumber\\
\qquad =(\tau \times \operatorname{Id})\circ (\phi \times \operatorname{Id})\circ (\operatorname{Id} \times \phi)\circ (\operatorname{Id} \times \tau).\label{eq24}
\end{gather}
$(\Omega,\rightarrow,\triangleright)$ is a dual CEDS if and only if
\begin{gather}
(\operatorname{Id} \times \phi)\circ (\phi\times \operatorname{Id})\circ (\operatorname{Id} \times \phi) =(\phi \times \operatorname{Id})\circ (\operatorname{Id} \times \tau)\circ (\phi\times \operatorname{Id}),\tag{\ref{eq23}}\\
(\phi \times \operatorname{Id})\circ (\tau \times \operatorname{Id})\circ (\operatorname{Id} \times \tau)\circ (\operatorname{Id} \times \phi) \nonumber\\
\qquad=(\operatorname{Id} \times \tau)\circ (\operatorname{Id} \times \phi)\circ (\phi \times \operatorname{Id})\circ (\tau \times \operatorname{Id}).\label{eq25}
\end{gather}\end{enumerate}\end{Lemma}

With this reformulation, the following result becomes immediate, as the
inversion of \eqref{eq23} gives \eqref{eq23} again and the inversion of \eqref{eq24} gives \eqref{eq25}.

\begin{Proposition}\label{prop2.7}
Let $(\Omega, \rightarrow,\triangleright)$ be a set with two binary operations. We shall say that $(\Omega,\rightarrow, \triangleright)$ is non-degenerate if the map $\phi$ of Definition~{\rm\ref{defi2.4}} is a bijection. If so, we put
\[\phi^{-1}\colon \ \begin{cases}
\Omega^2 \longrightarrow \Omega^2,\\
(\alpha,\beta) \longrightarrow (\alpha\curvearrowright \beta,\alpha \blacktriangleright \beta).
\end{cases}\]
Then $(\Omega,\rightarrow,\triangleright)$ is an EAS $($resp.\ a CEDS, a dual CEDS$)$ if and only if $(\Omega,\curvearrowright,\blacktriangleright)$ is an EAS $($resp.\ a dual CEDS, a CEDS$)$.
\end{Proposition}

\subsection{EAS of cardinality two}

Here is a classification of EAS of cardinality two, which we obtained by an exhaustive study of the $2^8$ possibilities of pairs of operations. The underlying set is $\Omega=\{X,Y\}$ and the products will be given by the pair of matrices
\begin{align*}
\begin{pmatrix}
X\rightarrow X&X\rightarrow Y\\Y\rightarrow X&Y\rightarrow Y
\end{pmatrix},\qquad
\begin{pmatrix}
X\triangleright X&X\triangleright Y\\Y\triangleright X&Y\triangleright Y
\end{pmatrix}.
\end{align*}
We shall use the two maps
\begin{align*}
\pi_X\colon \ \begin{cases}
\Omega\longrightarrow\Omega,\\
\alpha\longrightarrow X,
\end{cases}\qquad
\pi_Y\colon \ \begin{cases}
\Omega\longrightarrow\Omega,\\
\alpha\longrightarrow Y.
\end{cases}
\end{align*}
We respect the indexation of EDS of \cite{Foissy49}.
$$\renewcommand{\arraystretch}{1.1}
\begin{array}{|c|c|c|l|l|}
\hline \text{Case}&\rightarrow&\triangleright&\text{Description}&\text{Comments}\\
\hline \hline \mathbf{A1}&\begin{pmatrix}
X&X\\X&X
\end{pmatrix}&
\begin{pmatrix}
X&X\\X&X
\end{pmatrix}&\eas(\Omega,\rightarrow,\pi_X)&\begin{array}{@{}l@{}}\text{CEDS, dual CEDS,}\\
 \text{right part of \textbf{D1}}\end{array}\\
\hline \mathbf{A2}&\begin{pmatrix}
X&X\\X&X
\end{pmatrix}&
\begin{pmatrix}
X&X\\Y&Y
\end{pmatrix}&\eas(\Omega,\rightarrow)&\begin{array}{@{}l@{}}\text{CEDS, dual CEDS,}\\
\text{right part of \textbf{D2}}\end{array}\\
\hline \mathbf{C1} &\begin{pmatrix}
X&X\\X&Y
\end{pmatrix}&
\begin{pmatrix}
X&X\\X&X
\end{pmatrix}&\eas(\Omega,\rightarrow,\pi_X)&\text{CEDS, right part of \textbf{C4}}\\
\hline \mathbf{C3}&\begin{pmatrix}
X&X\\X&Y
\end{pmatrix}&
\begin{pmatrix}
X&X\\Y&Y
\end{pmatrix}&\eas(\mathbb{Z}/2\mathbb{Z},\times)&\text{CEDS, dual CEDS}\\
\hline \mathbf{C5} &\begin{pmatrix}
X&X\\X&Y
\end{pmatrix}&
\begin{pmatrix}
Y&Y\\Y&Y
\end{pmatrix}&\eas((\mathbb{Z}/2\mathbb{Z},\times),\pi_Y)&\text{CEDS, right part of \textbf{C2}}\\
\hline \mathbf{C6}&\begin{pmatrix}
X&X\\X&Y
\end{pmatrix}&
\begin{pmatrix}
X&X\\Y&X
\end{pmatrix}&&\\
\hline \mathbf{E1'}\text{--}\mathbf{E2'}&\begin{pmatrix}
X&X\\Y&Y
\end{pmatrix}&
\begin{pmatrix}
X&X\\X&X
\end{pmatrix}&\eas(\Omega,\rightarrow,\pi_X)&\text{right part of \textbf{E1} and \textbf{E2}}\\
\hline \mathbf{E3'}&\begin{pmatrix}
X&X\\Y&Y
\end{pmatrix}&
\begin{pmatrix}
X&X\\Y&Y
\end{pmatrix}&\eas(\Omega,\rightarrow)&\begin{array}{@{}l@{}}\text{dual CEDS},\\
\text{right part of \textbf{E3}}\end{array}\\
\hline\mathbf{ F1}&\begin{pmatrix}
X&Y\\X&Y
\end{pmatrix}&
\begin{pmatrix}
X&X\\X&X
\end{pmatrix}&\eas(\Omega,\rightarrow,\pi_X)&\begin{array}{@{}l@{}}\text{CEDS, dual CEDS},\\
\text{right part of \textbf{B1}, \textbf{F2}, \textbf{G1}}\\
\text{and \textbf{G2}}\end{array}\\
\hline \mathbf{F3}&\begin{pmatrix}
X&Y\\X&Y
\end{pmatrix}&
\begin{pmatrix}
X&X\\Y&Y
\end{pmatrix}&\eas(\Omega)&\begin{array}{@{}l@{}}\text{CEDS, dual CEDS},\\
\text{non-degenerate},\\
\text{right part of \textbf{B2} and \textbf{G3}}\end{array}\\
\hline \mathbf{F4}&\begin{pmatrix}
X&Y\\X&Y
\end{pmatrix}&
\begin{pmatrix}
X&Y\\Y&X
\end{pmatrix}&\eas'(\mathbb{Z}/2\mathbb{Z},+)&\begin{array}{@{}l@{}}\text{CEDS, dual CEDS},\\
\text{non-degenerate},\\
\text{right part of \textbf{F5}}\end{array}\\
\hline \mathbf{H1}&\begin{pmatrix}
X&Y\\Y&X
\end{pmatrix}&
\begin{pmatrix}
X&X\\X&X
\end{pmatrix}&\eas(\mathbb{Z}/2\mathbb{Z},+,\pi_X)&\text{CEDS}\\
\hline \mathbf{H2}&\begin{pmatrix}
X&Y\\Y&X
\end{pmatrix}&
\begin{pmatrix}
X&X\\Y&Y
\end{pmatrix}&\eas(\mathbb{Z}/2\mathbb{Z},+)&\begin{array}{@{}l@{}}\text{CEDS, dual CEDS},\\
\text{non-degenerate}\end{array}\\
\hline \end{array}
$$
For the cases \textbf{C3}, \textbf{C5}, \textbf{F4}, \textbf{H1} and \textbf{H2}, $\Omega$ is identified with $\mathbb{Z}/2\mathbb{Z}$, $X$ being $\overline{0}$ and $Y$ being $\overline{1}$.

\begin{Remark}
With similar methods, it is possible to prove that there are three non-degenerate EAS of cardinality 3 up to isomorphism: $\eas(\{1,2,3\})$, $\eas(\mathbb{Z}/3\mathbb{Z},+)$ and $\eas'(\mathbb{Z}/3\mathbb{Z},+)$. All of them are both CEDS and dual CEDS.
\end{Remark}


\section{Structure of non-degenerate finite CEDS}

We now turn to CEDS, and prove the structure Theorem~\ref{theo3.14} after several intermediate results.

\subsection{Preliminary results}

\begin{Lemma}\label{lem3.1}
let $\Omega$ be a finite non-degenerate EAS.
\begin{enumerate}\itemsep=0pt
\item[$(1)$] Let $\Omega'$ be a sub-EAS of $\Omega$. Then $\Omega'$ is non-degenerate.
\item[$(2)$] Let $\sim$ be an equivalence on $\Omega$, compatible with the EAS structure. Then the quotient EAS $\Omega/{\sim}$ is non-degenerate.
\end{enumerate}
\end{Lemma}

\begin{proof}
(1) By restriction, $\phi_{\Omega'}=(\phi_\Omega)_{\mid \Omega'^2}$ is injective. As $\Omega'$ is finite, it is a bijection. So $\Omega'$ is non degenerate.

(2) Let $\pi\colon\Omega\longrightarrow \Omega/{\sim}$ be the canonical surjection. Then $\phi_{\Omega/{\sim}}\circ \pi=(\pi\otimes \pi)\circ \phi_\Omega$. As $\phi$ is surjective, $\phi_{\Omega/{\sim}}$ is surjective.
As $\Omega/{\sim}$ is finite, it is a bijection. So $\Omega/{\sim}$ is non-degenerate.
 \end{proof}

\begin{Definition}
Let $(\Omega,\rightarrow,\triangleright)$ be an EAS. For any $\alpha \in \Omega$, we put
\begin{align*}
\phi_\alpha\colon\ \begin{cases}
\Omega \longrightarrow \Omega,\\
\beta \longrightarrow \alpha\rightarrow \beta,
\end{cases}\qquad
\psi_\alpha\colon\ \begin{cases}
\Omega \longrightarrow \Omega,\\
\beta \longrightarrow \beta\triangleright\alpha.
\end{cases}
\end{align*}
We shall say that $(\Omega,\rightarrow,\triangleright)$ is strongly non-degenerate if for any $\alpha\in \Omega$, $\phi_\alpha$ is bijective.
\end{Definition}

\begin{Remark}
As the product $\rightarrow$ is associative, for any $\alpha,\beta \in \Omega$,
$\phi_\alpha \circ \phi_\beta=\phi_{\alpha \rightarrow \beta}$.
\end{Remark}

\begin{Lemma} \label{lem3.3}
Let $(\Omega,*)$ be an associative semigroup. The following conditions are equivalent:
\begin{enumerate}\itemsep=0pt
\item[$(1)$] $\eas(\Omega,*)$ is non-degenerate.
\item[$(2)$] $\eas(\Omega,*^{\rm op})$ is strongly non-degenerate.
\item[$(3)$] $(\Omega,*^{\rm op})$ has the right inverse condition.
\end{enumerate}
\end{Lemma}

\begin{proof}
Let $\alpha$, $\beta$, $\gamma$, $\delta\in \Omega$. Then
\begin{align*}
\phi(\alpha,\beta)=(\gamma,\delta) \Longleftrightarrow \begin{cases}
\alpha*\beta=\gamma,\\
\alpha=\delta.
\end{cases}\end{align*}
So
\begin{align*}
\phi\text{ is bijective}&\Longleftrightarrow\forall (\gamma,\delta)\in \Omega^2,
\exists! \beta \in \Omega, \delta*\beta=\gamma\\
&\Longleftrightarrow \text{in $\eas(\Omega,*^{\rm op})$, }\forall \delta \in \Omega, \phi_{\delta}\text{ is bijective}\\
&\Longleftrightarrow\text{$(\Omega,*^{\rm op})$ has the right inverse condition}. \tag*{\qed}
\end{align*} \renewcommand{\qed}{}\end{proof}


\begin{Lemma}\label{lem3.4}
Let $(\Omega,\rightarrow,\triangleright)$ be a finite non-degenerate CEDS. Then it is strongly non-degenerate.
\end{Lemma}

\begin{proof}
Let $\alpha,\gamma,\gamma' \in \Omega$ such that $\phi_\alpha(\gamma)=\phi_\alpha(\gamma')$. In other words, $\alpha \rightarrow \gamma=\alpha\rightarrow \gamma'$. By \eqref{eq16},
\begin{align*}
\alpha \triangleright\gamma&=\alpha \triangleright (\alpha \rightarrow \gamma)=\alpha \triangleright (\alpha \rightarrow \gamma')=\alpha \triangleright \gamma'.
\end{align*}
Therefore, $\phi(\alpha,\gamma)=\phi(\alpha,\gamma')$. As $\phi$ is injective, $\gamma=\gamma'$, so $\phi_\alpha$ is injective. As $\Omega$ is finite, $\phi_\alpha$ is bijective.
\end{proof}

\begin{Lemma}\label{lem3.5}
Let $\Omega=(\Omega,\rightarrow,\triangleright)$ be a non-degenerate EAS, such that
$\forall \alpha,\beta \in \Omega$, $\alpha \rightarrow \beta =\beta$.
There exists a product $*$ on $\Omega$, making it a semigroup with the right inverse condition, such that~${\Omega=\eas'(\Omega,*)}$. For any $\beta\in \Omega$, $\psi_\beta$ is bijective and its inverse is
\begin{align*}
\phi'_\beta\colon\ \begin{cases}
\Omega \longrightarrow \Omega,\\
\alpha \longrightarrow \alpha*\beta.
\end{cases}
\end{align*}
Moreover, for any $\beta,\gamma \in \Omega$,
\begin{align}
\label{eq26}
\psi_\beta \circ \psi_\gamma =\psi_{\beta*\gamma},\qquad \psi_{\beta \triangleright \gamma} =\psi_\beta \circ \psi_\gamma^{-1}.
\end{align}\end{Lemma}

\begin{proof}
Note that for any $\alpha \in \Omega$, $\phi_\alpha=\operatorname{Id}_\Omega$. Let $\alpha,\beta,\gamma,\delta\in \Omega$. Then
\begin{align*}
\phi(\alpha,\beta)=(\gamma,\delta) \Longleftrightarrow \begin{cases}
\beta=\gamma,\\
\alpha\triangleright \beta=\delta.
\end{cases}\end{align*}
Hence,
\begin{align*}
\phi\text{ is bijective}&\Longleftrightarrow\forall (\gamma,\delta)\in \Omega^2, \exists! \alpha \in \Omega, \alpha \triangleright\gamma=\delta\Longleftrightarrow\forall\gamma \in \Omega, \psi_\gamma\text{ is bijective}.
\end{align*}
Putting $\phi^{-1}(\alpha,\beta)=(\alpha\curvearrowright \beta,\alpha \blacktriangleright \beta)$, by Proposition \ref{prop2.7} $(\Omega,\curvearrowright,\blacktriangleright)$ is an EAS, so $\curvearrowright$ is associative.
Moreover, $\phi^{-1}(\alpha,\beta)=(\alpha \curvearrowright \beta,\alpha)$, so $(\Omega,\curvearrowright,\blacktriangleleft)=\eas(\Omega,\curvearrowright)$.
By Lemma \ref{lem3.3}, if $*=\curvearrowright^{\rm op}$, then $*$ has the right inverse condition. Moreover, for any $\alpha,\beta \in \Omega$,
\begin{align*}
\phi^{-1}\circ \phi(\alpha,\beta)&=\phi^{-1}(\beta,\alpha \triangleright \beta)=((\alpha \triangleright \beta)*\beta,\beta)=(\alpha,\beta).
\end{align*}
Hence, the unique element $\gamma \in \Omega$ such that $\gamma*\beta=\alpha$ is $\alpha \triangleright \beta$: consequently, $\Omega=\eas'(\Omega,*)$. Moreover, for any $\alpha,\beta \in \Omega$,
$\phi'_\beta \circ \psi_\beta(\alpha)=(\alpha\triangleright \beta) *\beta=\alpha$.
So $\phi'_\beta\circ \psi_\beta=\operatorname{Id}_\Omega$. As $\psi_\beta$ is bijective, $\psi_\beta^{-1}=\phi'_\beta$.

Let $\beta,\gamma\in \Omega$. Then, for any $\alpha \in \Omega$,
$\phi'_\gamma \circ \phi'_\beta(\alpha)=\alpha *\beta*\gamma=\phi'_{\beta*\gamma}$.
So $\phi'_\gamma\circ \phi'_\beta=\phi'_{\beta*\gamma}$. Inverting, $\psi_\beta\circ \psi_\gamma=\psi_{\beta*\gamma}$. As a consequence,
$\psi_{\beta \triangleright \gamma}\circ \psi_\gamma=\psi_{(\beta \triangleright \gamma)*\gamma}=\psi_\beta$,
which induces the last formula. \end{proof}

\begin{Lemma}\label{lem3.6}
Let $\Omega=(\Omega,\rightarrow,\triangleright)$ be a non-degenerate EAS such that $\forall \alpha,\beta \in \Omega$,
$\alpha \rightarrow\beta=\beta$.
Then $\Omega_\psi=\{\psi_\alpha, \alpha\in \Omega\}$ is a subgroup of the group of permutations of $\Omega$.
\end{Lemma}

\begin{proof} Direct consequence of \eqref{eq26}. \end{proof}

\begin{Proposition}\label{prop3.7}
Let $\Omega=\eas'(\Omega,*)$, where $(\Omega,*)$ is a finite semigroup with the right inverse condition. We define an equivalence $\sim$ on $\Omega$ by $\alpha \sim \beta$ if $\psi_\alpha=\psi_\beta$. Then
\begin{enumerate}\itemsep=0pt
\item[$(1)$] $\sim$ is compatible with the EAS structure of $\Omega$. Therefore, $\Omega/{\sim}$ is an EAS.
\item[$(2)$] There exists a product $\star$ on $\Omega/{\sim}$, making it a group, such that $\Omega/{\sim}=\eas'(\Omega/{\sim},\star)$.
\item[$(3)$] There exists a sub-EAS $\Omega_0$ of $\Omega$, such that the restriction to $\Omega_0$ of the canonical surjection~${\pi\colon\Omega\longrightarrow \Omega/{\sim}}$ is an isomorphism.
\end{enumerate}\end{Proposition}

\begin{proof}
(1) Let $\alpha,\beta \in \Omega$, such that $\alpha \sim\beta$. Then $\psi_\alpha=\psi_\beta$. Let $\gamma \in \Omega$. Then $\alpha \rightarrow \gamma=\beta \rightarrow \gamma=\gamma$, and $\gamma \rightarrow \alpha=\alpha \sim \beta=\gamma \rightarrow \beta$.
As $\psi_\alpha=\psi_\beta$, $\gamma \triangleright \alpha=\gamma \triangleright \beta$. Moreover, by Lemma \ref{lem3.5},
\[\psi_{\alpha \triangleright \gamma}=\psi_\alpha \circ \psi_\gamma^{-1}=\psi_\beta\circ \psi_\gamma^{-1}=\psi_{\beta \triangleright \gamma},\]
so $\alpha \triangleright \gamma\sim \beta \triangleright \gamma$: $\sim$ is compatible with the EAS structure.

(2) By Lemma \ref{lem3.1}, $\Omega/{\sim}$ is non-degenerate. By Lemma \ref{lem3.5}, there exists a product $\star$ satisfying the right inverse condition, such that $\Omega/{\sim}=\eas'(\Omega/{\sim},\star)$. We consider the map
\[\psi\colon\ \begin{cases}
(\Omega/{\sim},\star) \longrightarrow (\mathfrak{S}_{\Omega/{\sim}},\circ),\\
\overline{\alpha} \longrightarrow \psi_{\overline{\alpha}}.
\end{cases}\]
By Lemma \ref{lem3.5}, this is a semigroup morphism. Let us prove that it is injective. We assume that $\psi_{\overline{\alpha}}=\psi_{\overline{\beta}}$.
In other words, for any $\gamma \in \Omega$, $\gamma \triangleright \alpha\sim \gamma \triangleright \beta$, or equivalently, $\psi_{\gamma \triangleright \alpha}=\psi_{\gamma \triangleright \beta}$. Moreover,
\[\psi_{\gamma \triangleright \alpha}=\psi_\gamma \circ \psi_\alpha^{-1}=\psi_{\gamma \triangleright \beta}=\psi_\gamma \circ \psi_\beta^{-1}.\]
As $\psi_\gamma$ is bijective, $\psi_\alpha=\psi_\beta$, so $\overline{\alpha}=\overline{\beta}$.

By Lemma \ref{lem3.6}, there exists $e\in \Omega/{\sim}$, such that $\psi_e=\operatorname{Id}_{\Omega/{\sim}}$. For any $\overline{\alpha}\in \Omega/{\sim}$,
\[\psi_{e\triangleright \overline{\alpha}}=\psi_e\circ \psi_{\overline{\alpha}}^{-1}=\psi_{\overline{\alpha}}^{-1},\]
so $\psi(\Omega/{\sim})$ is a subgroup of $\mathfrak{S}_{\Omega/{\sim}}$. Consequently, $(\Omega/{\sim},\star)$ is a group.

(3) By Lemma \ref{lem3.6}, there exists $\beta_0\in \Omega$ such that $\psi_{\beta_0}=\operatorname{Id}_\Omega$. We put
\[\Omega_0=\{\beta_0\triangleright \alpha,\, \alpha \in \Omega\}=\{\psi_\alpha(\beta_0),\, \alpha \in \Omega\}.\]
As the product $\rightarrow$ of $\Omega$ is trivial, this is a sub-semigroup of $(\Omega,\rightarrow)$. Let $\beta_0\triangleright \alpha$, $\beta_0\triangleright \beta \in \Omega_0$.
\[(\beta_0\triangleright \alpha)\triangleright (\beta_0\triangleright \beta)=\psi_{\beta_0\triangleright \gamma}\circ \psi_{\alpha}(\beta_0)=\psi_{(\beta_0\triangleright \gamma)*\alpha}(\beta_0)\in \Omega_0,\]
so $\Omega_0$ is a sub-EAS of $\Omega$.

Let us assume that $\beta_0\triangleright \alpha\sim \beta_0\triangleright \beta$. Then
\begin{align*}
\psi_{\beta_0\triangleright \alpha}=\psi_{\beta_0}\circ \psi_\alpha^{-1}=\psi_{\alpha}^{-1}
=\psi_{\beta_0\triangleright \beta}=\psi_{\beta_0}\circ \psi_\beta^{-1}=\psi_{\beta}^{-1},
\end{align*}
so $\psi_\alpha=\psi_\beta$. Hence, $\beta_0\triangleright\alpha=\beta_0\triangleright \beta$, which proves that \smash{$\pi_{\mid \Omega_0}$} is injective.
By Lemma \ref{lem3.6}, there exists $\beta \in \Omega$ such that $\psi_\beta=\psi_\alpha^{-1}$. We consider $\beta_0\triangleright \beta \in \Omega_0$. Then
$\psi_{\beta_0\triangleright \beta}=\psi_{\beta_0}\circ \psi_\beta^{-1}=\psi_\alpha$,
so $\beta_0\triangleright \beta\sim \alpha$. Hence, $\pi_{\mid \Omega_0}$ is surjective.
\end{proof}

\begin{Theorem} \label{theo3.8}
Let $\Omega=\eas'(\Omega,*)$, where $(\Omega,*)$ is a finite semigroup with the right inverse condition. There exists a group $(\Omega_0,\star)$ and a set $\Omega_1$ such that
$\Omega\approx \eas(\Omega_1)\times \eas'(\Omega_0,\star)$.
\end{Theorem}

\begin{proof}
We keep the notations of the proof of Proposition \ref{prop3.7}. As the sub-EAS $\Omega_0$ is isomorphic to $\Omega/{\sim}$, it is a group for the law $*$, and $\Omega_0=\eas'(\Omega_0,*)$.
Let $e$ be the unit of the group~$(\Omega/{\sim},\star)$. We consider
$\Omega_1=\{\alpha\in \Omega,\, \overline{\alpha}=e\}$.
Let us prove that $\Omega_1=\{\alpha \in \Omega,\, \psi_\alpha=\operatorname{Id}_\Omega\}$.

$\supseteq$: if $\psi_\alpha=\operatorname{Id}_\Omega$, for any $\beta \in \Omega$, $\overline{\beta} \star \overline{\alpha}^{*-1}=\overline{\psi_\alpha(\beta)}=\overline{\beta}$, so $\overline{\alpha}=e$ and $\alpha \in \Omega_1$.

$\subseteq$: if $\overline{\alpha}=e$, then for any $\beta \in \Omega$, $\overline{\beta \triangleright \alpha}=\overline{\beta}$, so $\beta \triangleright \alpha\sim \beta$: in other words, $\psi_{\beta \triangleright \alpha}=\psi_\beta$. Then
\[\psi_{\beta \triangleright \alpha}=\psi_\beta \circ \psi_{\alpha}^{-1}=\psi_\beta.\]
As $\psi_\beta$ is bijective, $\psi_\alpha=\operatorname{Id}_\Omega$.

Therefore, for any $\alpha \in \Omega$, for any $\beta \in \Omega_1$,
$\alpha \triangleright \beta=\psi_\beta(\alpha)=\alpha$.
As a consequence, $\Omega_1=\eas(\Omega_1)$. We consider the map
\[\theta\colon\ \begin{cases}
\Omega_1\times \Omega_0 \longrightarrow \Omega,\\
(\alpha,\beta) \longrightarrow \alpha*\beta.
\end{cases}\]
Let us prove that $\theta$ is injective. If $\theta(\alpha,\beta)=\theta(\alpha',\beta')$, in $\Omega/{\sim}$,
$\overline{\alpha}\star \overline{\beta}=\overline{\beta}=\overline{\alpha'}\star \overline{\beta'}=\overline{\beta'}$.
As $\pi_{\mid \Omega_0}$ is injective, $\beta=\beta'$. Because of the right inverse condition for $*$, $\alpha=\alpha'$.

Let us prove that $\theta$ is surjective. Let $\gamma \in \Omega$. There exists a unique $\beta\in \Omega_0$ such that $\psi_{\overline{\gamma}}=\psi_{\overline{\beta}}$.
We put $\alpha=\gamma \triangleright \beta$, so $\gamma=\alpha*\beta$. Moreover,
\smash{$\psi_{\overline{\alpha}}=\psi_{\overline{\gamma}}\circ \psi_{\overline{\beta}}^{-1}=\operatorname{Id}_{\Omega/{\sim}}$}.
So \smash{$e\triangleright \overline{\alpha}=\overline{\alpha}^{\star -1}=\psi_{\overline{\alpha}}(e)=e$}, and finally $\alpha\in \Omega_1$.

Let $(\alpha,\beta)$ and $(\alpha',\beta')\in \Omega_1 \times \Omega_0$. In $\Omega$, as $\alpha'\in \Omega_1$,
$\alpha*\beta*\alpha'*\beta'=\alpha*(\beta*\beta')$,
which implies that~${(\alpha*\beta)\triangleright (\alpha'*\beta')=\alpha* \bigl(\beta*\beta'^{*-1}\bigr)}$.
So $\theta$ is an isomorphism of EAS from $\eas(\Omega_1)\times \eas'(\Omega_0,*)$ to $\Omega$.
\end{proof}


\subsection{Non-degenerate finite CEDS}

\begin{Lemma}\label{lem3.9}
Let $(\Omega,\rightarrow,\triangleright)$ be a strongly non-degenerate finite EAS. Then $\Omega_\phi=(\{\phi_\alpha, \alpha\in \Omega\},\circ)$ is a group. The following map is a surjective morphism of semigroups:
\[\phi\colon \ \begin{cases}
(\Omega,\rightarrow) \longrightarrow \Omega_\phi,\\
\alpha \longrightarrow \phi_\alpha.
\end{cases}\]
\end{Lemma}

\begin{proof}
We already observed that for any $\alpha,\beta \in \Omega$, $\phi_\alpha\circ \phi_\beta=\phi_{\alpha\rightarrow \beta}$, so $\phi$ is a semigroup morphism.
By hypothesis, for any $\alpha \in \Omega$, $\phi_\alpha$ is a bijection, so belongs to the symmetric group~$\mathfrak{S}_\Omega$ of permutations of $\Omega$.
As $\Omega$ is finite, for any $\alpha \in \Omega$, there exists $n\geqslant 2$ such that~${\phi_\alpha^n=\operatorname{Id}_\Omega}$.
Then $\phi_{\alpha^{\rightarrow n}}=\operatorname{Id}_\Omega$, so $\Omega_\phi$ is a monoid. Putting $\beta=\alpha^{\rightarrow(n-1)}$, $\phi_\beta\circ \phi_\alpha=\phi_\alpha\circ \phi_\beta=\operatorname{Id}_\Omega$, so $\Omega_\phi$ is a group.
\end{proof}

\begin{Proposition}\label{prop3.10}
Let $\Omega=(\Omega,\rightarrow,\triangleright)$ be a finite non-degenerate EAS, such that for any $\alpha \in \Omega$, $\phi_\alpha$ is a bijection. We put
$
\Omega^\rightarrow=\{\alpha\in \Omega, \phi_\alpha=\operatorname{Id}_\Omega\}$, $\Omega^\triangleright =\{\beta \in \Omega, \psi_\beta=\operatorname{Id}_\Omega\}$.
Then
\begin{enumerate}\itemsep=0pt
\item[$(1)$] $\Omega^\rightarrow$ is a non-degenerate sub-EAS of $\Omega$.
\item[$(2)$] If $\Omega^\triangleright$ is nonempty, it is a non-degenerate sub-EAS of $\Omega$.
\item[$(3)$] If $\Omega^\triangleright$ is nonempty, then $\Omega^\triangleright \cap \Omega^\rightarrow$ is nonempty.
\item[$(4)$] If $\Omega$ is a CEDS, $\Omega^\triangleright$ is nonempty.
\end{enumerate}\end{Proposition}

\begin{proof}
(1) Recall that for any $\alpha,\beta \in \Omega$, $\phi_\alpha\circ \phi_\beta=\phi_{\alpha \rightarrow \beta}$. This easily implies that $\Omega^\rightarrow$ is stable under $\rightarrow$.
By Lemma \ref{lem3.9}, there exists $\alpha \in \Omega$, such that $\phi_\alpha=\operatorname{Id}_\Omega$, so $\Omega^\rightarrow$ is nonempty.

Let $\alpha,\beta \in \Omega^\rightarrow$. Let us consider $\gamma \in \Omega$. As $\phi$ is bijective, there exist $\beta',\gamma'\in \Omega$ such that
\[(\beta'\rightarrow \gamma',\beta'\triangleright \gamma')=(\beta,\gamma).\]
Then
\begin{align*}
\phi_{\alpha \triangleright \beta}(\gamma)&=(\alpha\triangleright \beta)\rightarrow \gamma=\bigl(\alpha \triangleright\bigl(\beta'\rightarrow \gamma'\bigr)\bigr)\rightarrow \bigl(\beta' \triangleright \gamma'\bigr)=\beta' \triangleright \gamma'=\gamma.
\end{align*}
So $\phi_{\alpha \triangleright \beta}=\operatorname{Id}_\Omega$ and $\alpha \triangleright \beta \in \Omega^\rightarrow$. By Lemma \ref{lem3.1}, $\Omega^\rightarrow$ is a non-degenerate sub-EAS.

(2) Let $\beta,\gamma \in \Omega^\triangleright$. As $\psi_\gamma=\operatorname{Id}_\Omega$, $\beta\triangleright \gamma=\beta \in \Omega^\triangleright$. For any $\alpha \in \Omega$, by \eqref{eq12} and \eqref{eq13},
\begin{gather*}
(\alpha \triangleright (\beta\rightarrow \gamma))\rightarrow (\beta \triangleright \gamma) =(\alpha \rightarrow \beta)\triangleright \gamma=\alpha \rightarrow \beta,\\
(\alpha \triangleright (\beta\rightarrow \gamma))\triangleright (\beta \triangleright \gamma) =\alpha \triangleright \beta.
\end{gather*}
So $\phi(\alpha \triangleright (\beta\rightarrow \gamma),\beta \triangleright \gamma)=\phi(\alpha,\beta)$.
As $\phi$ is injective, $\alpha \triangleright (\beta\rightarrow \gamma)=\alpha$, so $\psi_{\beta \rightarrow \gamma}=\operatorname{Id}_\Omega$ and~${\beta \rightarrow \gamma\in \Omega^\triangleright}$.
If $\Omega^\triangleright$ is nonempty, by Lemma \ref{lem3.1}, it is non-degenerate.

(3) Let us take $\alpha \in \Omega^\triangleright$. The permutation $\phi_\alpha$ is of finite order as $\Omega$ is finite, so there exists~${n\geqslant 2}$, such that $\phi_{\alpha}^n=\phi_{\alpha^{\rightarrow n}}=\operatorname{Id}_\Omega$.
Putting $\beta=\alpha^{\rightarrow n}$, then $\beta \in \Omega^{\triangleright}$ (as it is a sub-CEDS) and $\beta \in \Omega^\rightarrow$ as $\phi_\beta=\operatorname{Id}_\Omega$.

(4) Let us consider the EAS associated to the inverse of $\phi$ (see Proposition \ref{prop2.7}), which we denote by $(\Omega,\curvearrowright,\blacktriangleright)$.
By the first point, there exists $\alpha \in \Omega$ such that for any $\beta \in \Omega$, $\alpha \curvearrowright \beta=\beta$. In other words, for any $\beta \in \Omega$,
$\phi^{-1}(\alpha,\beta)=(\beta,\alpha \blacktriangleright \beta)$.
This implies that $\phi_\beta(\alpha \blacktriangleright \beta)=\alpha$. The inverse of the bijection $\phi_\beta$ is the map
\[\phi'_\beta\colon \ \begin{cases}
\Omega \longrightarrow \Omega,\\
\alpha \longrightarrow \alpha \blacktriangleright \beta.
\end{cases}\]
As $\Omega$ is finite, there exists $\beta'\in \Omega$ such that $\phi_\beta^{-1}=\phi_{\beta'}$. Hence,
\[\beta=\beta \triangleright (\alpha \blacktriangleright \beta)=\beta\triangleright \bigl(\beta'\rightarrow \alpha\bigr)=\beta \triangleright \alpha,\]
by \eqref{eq16}. So $\alpha \in \Omega^\triangleright$. \end{proof}

\begin{Proposition}\label{prop3.11}
Let $(\Omega,\rightarrow,\triangleright)$ be a finite non-degenerate CEDS.
\begin{enumerate}\itemsep=0pt
\item[$(1)$] We define an equivalence $\equiv$ on $\Omega$ by
$\beta \equiv \beta'$ if $\exists \alpha \in \Omega$, $\beta' =\alpha\rightarrow \beta$.
This equivalence is compatible with the CEDS structure. Therefore, $\Omega/\equiv$ is a CEDS.
\item[$(2)$] The restriction to $\Omega^\rightarrow$ of the canonical surjection $\pi\colon\Omega\longrightarrow \Omega/\equiv$ is an isomorphism.
\item[$(3)$] $\Omega=\Omega^\triangleright\rightarrow \Omega^\rightarrow$.
\end{enumerate}\end{Proposition}

\begin{proof}
(1) The relation $\equiv$ can be reformulated as: there exists $\phi_\alpha\in \Omega_\phi$, such that $\phi_\alpha(\beta)=\beta'$. By Lemmas \ref{lem3.4} and \ref{lem3.9}, $\Omega_\phi$ is a group.
This easily implies that $\equiv$ is an equivalence: its classes are the orbits of the action of the group $\Omega_\phi$ over $\Omega$.

Let us assume that $\beta \equiv \beta'$: we put $\beta'=\alpha\rightarrow \beta$. Let $\gamma \in \Omega$. Then $\gamma \rightarrow \beta\equiv \beta\equiv \beta'\equiv \gamma \rightarrow \beta'$ by definition of $\equiv$.
Moreover, $\beta' \rightarrow \gamma=\alpha \rightarrow (\beta \rightarrow \gamma)$, so $\beta'\rightarrow \gamma\equiv \beta \rightarrow \gamma$. By \eqref{eq13},
\begin{gather*}
\beta'\triangleright \gamma=(\alpha \rightarrow \gamma)\triangleright \gamma=(\alpha \triangleright (\beta \rightarrow \gamma))\rightarrow (\beta \triangleright \gamma)\equiv \beta \triangleright \gamma,\\
\gamma \triangleright \beta'=\gamma \triangleright(\alpha \rightarrow \beta)=\gamma \triangleright \beta.
\end{gather*}
So $\equiv$ is compatible with the CEDS structure.

(2) Let $\alpha \in \Omega$. As $\phi_\alpha$ is bijective, there exists a unique $\beta \in \Omega$ such that $\alpha \rightarrow \beta=\alpha$. Then
$\phi_\alpha=\phi_{\alpha \rightarrow \beta}=\phi_\alpha\circ \phi_\beta$.
As $\phi_\alpha$ is bijective, $\phi_\beta=\operatorname{Id}_\Omega$, so $\beta \in \Omega^\rightarrow$ and $\alpha \equiv \beta$. This proves that~$\pi_{\mid \Omega^\rightarrow}$ is surjective.

Let $\beta,\beta'\in \Omega^\rightarrow$, such that $\beta\equiv \beta'$. There exists $\alpha\in \Omega$ such that $\alpha\rightarrow\beta=\beta'$. Then
\[\operatorname{Id}_\Omega=\phi_{\beta'}=\phi_{\alpha}\circ \phi_{\beta}=\phi_\alpha,\]
so $\phi_\alpha=\operatorname{Id}_\Omega$. We deduce that $\beta'=\phi_\alpha(\beta)=\beta$. Hence, $\pi_{\mid \Omega^\rightarrow}$ is injective.

(3) By Proposition \ref{prop3.10}, there exists $\beta_0\in \Omega^\triangleright \cap \Omega^\rightarrow$.
Let $\beta \in \Omega$. As $\pi_{\mid \Omega^\rightarrow}$ is bijective, there exists $\beta_1\in \Omega^\rightarrow$, such that $\beta\equiv \beta_1$. We put $\beta=\alpha \rightarrow \beta_1$. As $\beta_0\in \Omega^\rightarrow$,
$\beta=\alpha \rightarrow {\beta_0\rightarrow \beta_1}$. Moreover, for any $\gamma \in \Omega$, as $\beta_0\in \Omega^\triangleright$, by \eqref{eq16},
$\gamma \triangleright (\alpha \rightarrow \beta_0)=\gamma \triangleright \beta_0=\gamma$,
so $\alpha \rightarrow \beta_0\in \Omega^\triangleright$. \end{proof}

\begin{Proposition}\label{prop3.12}
Let $(\Omega,\rightarrow,\triangleright)$ be a finite non-degenerate CEDS. We define an equivalence on $\Omega^\triangleright$ by
\[\alpha' \sim \alpha'' \Longleftrightarrow \exists \alpha \in \Omega^\rightarrow, \ \alpha''=\alpha'\rightarrow \alpha.\]
\begin{enumerate}\itemsep=0pt
\item[$(1)$] This equivalence is compatible with the CEDS structure, and $\Omega'=\Omega^\triangleright/{\sim}$ is a non-degen\-er\-ate CEDS. Moreover, $(\Omega',\rightarrow)$ is an abelian group and $\Omega'=\eas(\Omega',\rightarrow)$.
\item[$(2)$] The following map is a semigroup isomorphism:
\[\theta\colon \ \begin{cases}
(\Omega' \times \Omega^\rightarrow,\rightarrow) \longrightarrow (\Omega,\rightarrow),\\
(\overline{\alpha},\beta) \longrightarrow \alpha \rightarrow \beta.
\end{cases}\] \end{enumerate} \end{Proposition}

\begin{proof} We firstly introduce an auxiliary map, defined by
\[\Theta\colon \ \begin{cases}
\Omega^\triangleright \times \Omega^\rightarrow \longrightarrow \Omega,\\
(\alpha,\beta) \longrightarrow \alpha \rightarrow \beta.
\end{cases}\]
By Proposition \ref{prop3.11}, it is surjective. Let us prove that $\Theta(\alpha',\beta')=\Theta(\alpha'',\beta'')$ if and only if $\alpha'\sim \alpha''$ and $\beta'=\beta''$.

Let us assume that $\Theta(\alpha',\beta')=\Theta(\alpha'',\beta'')$. As $\phi_{\alpha'}$ is bijective, there exists $\alpha \in \Omega$, $\alpha''=\alpha'\rightarrow \alpha$.
As $\alpha'\rightarrow \beta'=\alpha''\rightarrow \beta''$ and $\beta',\beta''\in \Omega^\rightarrow$,
\[\phi_{\alpha'}=\phi_{\alpha'}\circ \phi_{\beta'}=\phi_{\alpha'\rightarrow \beta'}=\phi_{\alpha''\rightarrow \beta''}=\phi_{\alpha''}.\]
Hence,
$\phi_{\alpha''}=\phi_{\alpha'}\circ \phi_{\alpha}=\phi_{\alpha'}$.
As $\phi_{\alpha'}$ is bijective, $\phi_\alpha=\operatorname{Id}_\Omega$, so $\alpha \in \Omega^\rightarrow$: we obtain that $\alpha'\equiv \alpha''$. As $\phi_{\alpha'}=\phi_{\alpha''}$, $\phi_{\alpha'}(\beta')=\phi_{\alpha''}(\beta'')=\phi_\alpha'(\beta'')$.
As $\phi_{\alpha'}$ is injective, $\beta=\beta'$. Conversely, if $\alpha \in \Omega^\rightarrow$, $\alpha'\rightarrow \alpha \rightarrow\beta'=\alpha'\rightarrow \beta'$.

As a consequence, $\sim$ is indeed an equivalence, $\theta$ is well-defined and is a bijection. It remains to show that $\sim$ is compatible with the CEDS structure of $\Omega^\triangleright$.
Let $\alpha',\alpha''\in \Omega^\triangleright$, such that $\alpha' \sim \alpha''$. We put $\alpha''=\alpha'\rightarrow \alpha$, with $\alpha \in \Omega^\rightarrow$. Let $\beta \in \Omega^\triangleright$. Then
\begin{align*}
\alpha'\triangleright \beta=\alpha'\sim \alpha''=\alpha''\triangleright \beta,\qquad\beta \triangleright \alpha'=\beta=\beta \triangleright \alpha''.
\end{align*}
Moreover,
\begin{align*}
\alpha''\rightarrow \beta=\alpha''\rightarrow \alpha \rightarrow \beta=\alpha''\rightarrow \beta,\qquad\beta \rightarrow \alpha''=\beta \rightarrow \alpha'\rightarrow \alpha\sim \beta \rightarrow \alpha'.
\end{align*}
Therefore, $\Omega^\triangleright/{\sim}$ is a CEDS. By Lemma \ref{lem3.1}, it is non-degenerate.

Let $\alpha,\alpha'\in \Omega^\triangleright$ and $\beta,\beta'\in \Omega^\rightarrow$. As $\beta\in \Omega^\rightarrow$,
\begin{align*}
\theta(\overline{\alpha},\beta)\rightarrow \theta\bigl(\overline{\alpha'},\beta'\bigr)&=\alpha \rightarrow \beta \rightarrow \alpha' \rightarrow \beta'=\alpha \rightarrow \alpha' \rightarrow \beta'\\
&=\alpha \rightarrow \alpha' \rightarrow \beta \rightarrow\beta'=\theta\bigl(\overline{\alpha}\rightarrow \overline{\alpha'},\beta \rightarrow \beta'\bigr).
\end{align*}
So $\theta$ is an isomorphism for the products $\rightarrow$.

Let us now study the CEDS $\Omega'$. By definition of $\Omega^\triangleright$, for any $\overline{\alpha}$, $\overline{\beta}\in \Omega'$, $\overline{\alpha}\triangleright \overline{\beta}=\overline{\alpha}$, so $\Omega'=\eas'(\Omega',\rightarrow)$.
By Proposition \ref{prop3.10}, $\Omega^\rightarrow$ is nonempty. Let us prove that
\[\Omega'^\rightarrow=\big\{\overline{\alpha}, \alpha\in \Omega^\triangleright\cap \Omega^\rightarrow\big\}.\]
$\supseteq$ is obvious. Let us take $\overline{\alpha}\in \Omega'^\rightarrow$. Then, for any $\beta \in \Omega^\triangleright$, $\overline{\alpha}\rightarrow\overline{\beta}=\overline{\beta}$:
there exists $\gamma \in \Omega^\rightarrow$, $\alpha \rightarrow \beta=\beta \rightarrow \gamma$. Therefore,
$\phi_\alpha\circ \phi_\beta=\phi_\beta \circ \phi_\gamma=\phi_\beta$,
as $\phi_\gamma=\operatorname{Id}_\Omega$. As $\phi_\beta$ is a bijection, $\phi_\alpha=\operatorname{Id}_\Omega$, so $\alpha \in \Omega^\rightarrow$. Let $\alpha,\beta \in \Omega^\triangleright\cap \Omega^\rightarrow$.
As $\phi_\alpha$ is bijective, there exists $\beta'\in \Omega$, $\alpha \rightarrow \beta'=\beta$. Then
\[ \operatorname{Id}_\Omega=\phi_\beta=\phi_\alpha\circ \phi_{\beta'}=\phi_{\beta'},\]
so $\beta'\in \Omega^\rightarrow$: we proved that $\alpha \sim \beta$. As a conclusion, there exists a unique $\overline{e}\in \Omega'$, such that for any $\overline{\alpha}\in \Omega'$, $\overline{e}\rightarrow \overline{\alpha}=\overline{\alpha}$.

Let us choose $\overline{\alpha}\in \Omega'$. As $\phi_\alpha$ is bijective, there exists $\overline{e'}\in \Omega'$ such that $\overline{\alpha}\rightarrow \overline{e'}=\overline{\alpha}$.
Let~${\overline{\beta}\in \Omega'}$. Then $\overline{\alpha}\rightarrow \overline{e'}\rightarrow \overline{\beta}=\overline{\alpha}\rightarrow \overline{\beta}$:
in other words, $\alpha \rightarrow e'\rightarrow \beta\sim \alpha \rightarrow \beta$, and there exists~${\gamma \in \Omega^\rightarrow}$, such that $\alpha \rightarrow e'\rightarrow \beta=\alpha \rightarrow \beta \rightarrow \gamma$.
As $\phi_\alpha$ is injective, $e'\rightarrow \beta=\beta \rightarrow \gamma\sim \beta$, so~${\overline{e'}\rightarrow \overline{\beta}=\overline{\beta}}$ for any $\overline{\beta}\in \Omega'$.
By unicity of $\overline{e}$, $\overline{e'}=\overline{e}$: for any $\overline{\alpha}\in \Omega'$, $\overline{\alpha}\rightarrow\overline{e}=\overline{\alpha}$, so $\overline{e}$ is a unit of $(\Omega',\rightarrow)$.
By~\eqref{eq15}, for $\gamma=\overline{e}$, we deduce that $(\Omega',\rightarrow)$ is an abelian monoid.
Let $\overline{\alpha}\in \Omega'$. As~${\phi_\alpha}$ is surjective, there exists $\overline{\alpha}'\in \Omega'$ such that $\overline{\alpha}\rightarrow \overline{\alpha}'=\overline{e}$. So $(\Omega',\rightarrow)$ is a group. \end{proof}

{\samepage
\begin{Proposition}
Let $(\Omega,*)$ be an associative semigroup such that for any $\alpha,\beta,\gamma \in \Omega$,
\[\alpha*\beta * \gamma=\beta *\alpha*\gamma.\]
Let $(\Omega',\rightarrow,\triangleright)$ be a CEDS, and $\prec\colon\Omega \times \Omega'\longrightarrow \Omega$ be a map such that for any $\alpha,\beta \in \Omega$, for any~${\beta',\gamma'\in \Omega'}$,
\begin{gather}
\label{eq27} \alpha \prec \bigl(\beta'\rightarrow \gamma'\bigr)=\alpha \prec \gamma',\\
\label{eq28} \bigl(\alpha *\beta\bigr)\prec \gamma' =\bigl(\alpha \prec \gamma'\bigr)*\bigl(\beta \prec \gamma'\bigr),\\
\label{eq29} \bigl(\alpha\prec \gamma'\bigr)\prec \bigl(\beta'\triangleright \gamma'\bigr)=\alpha \prec \beta',
\end{gather}
we define two products $\rightarrow$ and $\triangleright$ on $\Omega\times \Omega'$ in the following way: for any $(\alpha,\alpha')$, $(\beta,\beta')\in \Omega \times \Omega'$,
\begin{gather*}
\bigl(\alpha,\alpha'\bigr)\rightarrow \bigl(\beta,\beta'\bigr)=\bigl(\alpha*\beta,\alpha'\rightarrow \beta'\bigr),\qquad
\bigl(\alpha,\alpha'\bigr)\triangleright \bigl(\beta,\beta'\bigr)=\bigl(\alpha \prec \beta',\alpha'\triangleright \beta'\bigr).
\end{gather*}
Then $(\Omega\times \Omega,'\rightarrow,\triangleright)$ is a CEDS, which we denote by $\Omega \rtimes_\prec \Omega'$.
\end{Proposition}

\begin{proof}
Direct verifications.
\end{proof}

}

\begin{Remark}
If for any $(\alpha,\alpha')\in \Omega\times \Omega'$, $\alpha \prec \alpha'=\alpha$, we recover the direct product $\Omega \times \Omega'$ of EAS.
\end{Remark}


\begin{Theorem}\label{theo3.14}
Let $\Omega$ be a finite non-degenerate CEDS. There exist an abelian group $(\Omega_1,*)$, a group $(\Omega_2,\star)$, a left action $\succ\colon\Omega_2\times \Omega_1\longrightarrow \Omega_1$ of $(\Omega_2,\star)$ on $(\Omega_1,*)$
by group automorphisms, and a nonempty set $\Omega_3$ such that $\Omega$ is of the form
\[(\eas(\Omega_1,*) \rtimes_\succ \eas'(\Omega_2,\star))\times \eas(\Omega_3),\]
with the products given by
\begin{gather*}
(\alpha_1,\alpha_2,\alpha_3)\rightarrow(\beta_1,\beta_2,\beta_3)=(\alpha_1*\beta_1,\beta_2,\beta_3),\\
(\alpha_1,\alpha_2,\alpha_3)\triangleright(\beta_1,\beta_2,\beta_3)=\bigl(\beta_2\succ \alpha_1,\alpha_2\star \beta_2^{\star-1},\alpha_3\bigr).
\end{gather*}
\end{Theorem}



\begin{proof}
Let us consider the map $\theta$ of Proposition \ref{prop3.12}. For any $\alpha,\alpha'\in \Omega^\triangleright$, and for any $\beta,\beta'\in \Omega^\rightarrow$, by \eqref{eq16} and \eqref{eq17},
\begin{align*}
(\alpha \rightarrow \beta) \triangleright \bigl(\alpha'\rightarrow \beta'\bigr)&=(\alpha \rightarrow \beta)\triangleright \beta'=\bigl(\alpha \triangleright \beta'\bigr)\rightarrow \bigl(\beta \triangleright \beta'\bigr).
\end{align*}
Let $\beta_0\in \Omega^\triangleright \cap \Omega^\rightarrow$. Then, as $\beta_0\in \Omega^\rightarrow$,
\[(\alpha \rightarrow \beta) \triangleright \bigl(\alpha'\rightarrow \beta'\bigr)=\underbrace{\bigl(\alpha \triangleright \beta'\bigr)\rightarrow \beta_0}_{=\gamma_1}\rightarrow\underbrace{\bigl(\beta \triangleright \beta'\bigr)}_{=\gamma_2}.\]
Obviously, $\gamma_2\in \Omega^\rightarrow$. For any $\gamma \in \Omega$, by \eqref{eq16},
$\gamma \triangleright \gamma_1=\gamma \triangleright \beta_0=\gamma$,
so $\gamma_1\in \Omega^\triangleright$. We then put, for any $\overline{\alpha}\in \Omega'$, and for any $\beta\in \Omega^\rightarrow$,
$\overline{\alpha}\prec \beta=\overline{\alpha\triangleright \beta}\rightarrow \overline{\beta_0}$.
Then, for any $\overline{\alpha}, \overline{\alpha'}\in \Omega'$, and for any~${\beta,\gamma \in \Omega^\rightarrow}$,
$\theta(\overline{\alpha},\beta)\triangleright \theta\bigl(\overline{\alpha'},\beta'\bigr)=\theta\bigl(\overline{\alpha}\prec \beta',\beta \triangleright \beta'\bigr)$.
Then
\begin{align*}
\overline{\alpha}\prec(\beta \rightarrow \gamma)&=\overline{\alpha \triangleright(\beta \rightarrow \gamma)}\rightarrow \overline{\beta_0}=\overline{\alpha \triangleright \gamma}\rightarrow \overline{\beta_0}=\overline{\alpha}\prec \gamma,
\end{align*}
which proves \eqref{eq27}. As $\beta_0\in \Omega^\rightarrow$,
\begin{align*}
(\overline{\alpha}\rightarrow \overline{\beta})\prec \gamma&=\overline{(\alpha \rightarrow \beta) \triangleright \gamma}\rightarrow \overline{\beta_0}=\overline{\alpha \triangleright \gamma}\rightarrow\overline{\beta \triangleright \gamma}\rightarrow \overline{\beta_0}\\
&=\overline{\alpha \triangleright \gamma}\rightarrow \overline{\beta_0}\rightarrow\overline{\beta \triangleright \gamma}\rightarrow \overline{\beta_0}=(\overline{\alpha}\prec \gamma)\rightarrow (\overline{\beta}\prec \gamma),
\end{align*}
which proves \eqref{eq28}. Then
\begin{align*}
(\overline{\alpha}\prec \gamma)\prec (\beta \triangleright\gamma)&=\overline{(\alpha \triangleright \gamma)\triangleright (\beta \triangleright \gamma)}\rightarrow \overline{\beta_0\triangleright(\beta \triangleright \gamma)}\rightarrow \overline{\beta_0}=\overline{\alpha \triangleright \beta}\rightarrow \overline{\beta_0\triangleright(\beta \triangleright \gamma)}\rightarrow \overline{\beta_0}\\
&=\overline{\alpha \triangleright \beta}\rightarrow \overline{\beta_0}=\overline{\alpha}\prec \beta,
\end{align*}
which proves \eqref{eq29}. For the last equality, we used that $\beta_0\triangleright(\beta \triangleright \gamma)\in \Omega^\rightarrow$, as $\beta$, $\beta_0$ and $\gamma$ belong to $\Omega^\rightarrow$.

We finally obtain that $\theta$ is an isomorphism between $\Omega'\rtimes_\prec \Omega^\rightarrow$ and $\Omega$. We put $\Omega'=\eas(\Omega_1,*)$.
From Theorem~\ref{theo3.8}, we obtain a decomposition of $\Omega^\rightarrow$ of the form $\eas'(\Omega_2,\star)\times \eas(\Omega_3)$.
The map $\prec\colon\Omega_1\times \Omega_2\times \Omega_3\longrightarrow \Omega_1$ satisfies \eqref{eq27}--\eqref{eq29}. In this particular case, \eqref{eq27} becomes trivial, and \eqref{eq28}, \eqref{eq29} can be reformulated in this way:
for any $\alpha_1,\beta_1\in \Omega_1$, $\beta_2,\gamma_2\in \Omega_2$, $\beta_3,\gamma_3\in \Omega_3$,
\begin{gather*}
(\alpha_1*\beta_1)\prec(\gamma_2,\gamma_3)=(\alpha_1\prec (\gamma_2,\gamma_3))*(\beta_1\prec (\gamma_2,\gamma_3)),\\
(\alpha_1\prec (\gamma_2,\gamma_3))\prec(\beta_2,\beta_3)=\alpha_1\prec (\beta_2\star \gamma_2,\beta_3).
\end{gather*}
The products of $\Omega$ are given in this way: for any $\alpha_i,\beta_i\in \Omega_i$, with $1\leqslant i\leqslant 3$,
\begin{gather*}
\begin{split}
& (\alpha_1,\alpha_2,\alpha_3)\rightarrow(\beta_1,\beta_2,\beta_3)=(\alpha_1*\beta_1,\beta_2,\beta_3),\\
& (\alpha_1,\alpha_2,\alpha_3)\triangleright(\beta_1,\beta_2,\beta_3)=\bigl(\alpha_1\prec (\beta_2,\beta_3),\alpha_2\star \beta_2^{\star-1},\alpha_3\bigr).
\end{split}
\end{gather*}
For any $(\beta_2,\beta_3)\in \Omega_2\times \Omega_3$, we consider
\[\psi^\prec_{\beta_2,\beta_3}\colon \ \begin{cases}
\Omega_1 \longrightarrow \Omega_1,\\
\alpha_1 \longrightarrow \alpha_1\prec (\beta_2,\beta_3).
\end{cases}\]
As $\Omega$ is non-degenerate, necessarily $\psi_{\beta_2,\beta_3}$ is injective. As $\Omega$ is finite, $\psi_{\beta_2,\beta_3}$ is a bijection.
Moreover, by \eqref{eq30}, for any $(\beta_2,\beta_3)$, $(\gamma_2,\gamma_3)\in \Omega_2\times \Omega_3$,
$\psi^\prec_{\beta_2,\beta_3}\circ \psi^\prec_{\gamma_2,\gamma_3}=\psi^\prec_{\beta_2\star \gamma_2,\beta_3}$.
For $\beta_2=\gamma_2$ being the unit $e$ of $\Omega_2$ and $\beta_3=\gamma_3$, we obtain that \smash{$\bigl(\psi^\prec_{e,\beta_3}\bigr)^2=\psi^\prec_{e,\beta_3}$}.
As it is a~bijection, $\psi^\prec_{e,\beta_3}=\operatorname{Id}_{\Omega_1}$ for any $\beta_3\in \Omega_3$. Hence,
\[\psi^\prec_{e,\beta_3}\circ \psi^\prec_{\gamma_2,\gamma_3}=\psi^\prec_{\gamma_2,\gamma_3}=\psi^\prec_{\gamma_2,\beta_3},\]
so $\psi_{\beta_2,\beta_3}$ does not depend on $\beta_3$. We denote this map by $\psi_{\beta_2}$. Note that we proved that $\psi_{\beta_{e_{\Omega_2}}}=\operatorname{Id}_{\Omega_1}$.
We put, for any $\alpha_1\in \Omega_1$, $\beta_2\in \Omega_2$, $\alpha_1\prec \beta_2=\psi^\prec_{\beta_2}(\alpha_1)$. We finally obtain that the products in $\Omega$ are given by
\begin{gather*}
(\alpha_1,\alpha_2,\alpha_3)\rightarrow(\beta_1,\beta_2,\beta_3)=(\alpha_1*\beta_1,\beta_2,\beta_3),\\
(\alpha_1,\alpha_2,\alpha_3)\triangleright(\beta_1,\beta_2,\beta_3)=\bigl(\alpha_1\prec\beta_2,\alpha_2\star \beta_2^{\star-1},\alpha_3\bigr).
\end{gather*}
So $\Omega=(\eas(\Omega_1,*)\rtimes_\prec \eas'(\Omega_2,\star))\times \eas(\Omega_3)$.

In the particular case of $\eas(\Omega_1,*) \rtimes_\prec \eas'(\Omega_2,\star)$, \eqref{eq27} is trivial, and \eqref{eq28}, \eqref{eq29} can be reformulated in this way: for any $\alpha_1,\beta_1\in \Omega_1$, $\beta_2,\gamma_2\in \Omega_2$,
\begin{gather*}
(\alpha_1*\beta_1)\prec\gamma_2 =(\alpha_1\prec \gamma_2)*(\beta_1\prec \gamma_2),\qquad
(\alpha_1\prec \gamma_2)\prec \beta_2 =\alpha_1\prec (\beta_2\star \gamma_2).
\end{gather*}
As $\psi_{e_{\Omega_2}}=\operatorname{Id}_{\Omega_1}$, the following map is a left action of $(\Omega_2,\star)$ on $(\Omega_1,*)$ by group automorphisms:
\[\succ\colon \ \begin{cases}
\Omega_2\times \Omega_1 \longrightarrow \Omega_1,\\
(\beta_2,\alpha_1) \longrightarrow \beta_2\succ \alpha_1=\alpha_1\prec \beta_2.
\end{cases}\]
The formulas for the products in $\Omega$ are then immediate.
\end{proof}


\begin{Remark}
Consequently, we have a semi-direct product of groups $(\Omega_1,*)\rtimes_\succ (\Omega_2,\star)$.
\end{Remark}

Inverting the corresponding maps $\phi$, we obtain the following corollary.

\begin{Corollary}\label{cor3.15}
Let $\Omega$ be a finite non-degenerate dual CEDS. There exists an abelian group $(\Omega_1,*)$, a group $(\Omega_2,\star)$, a right action $\prec\colon\Omega_1\times \Omega_2\longrightarrow \Omega_1$ of $(\Omega_2,\star)$ on $(\Omega_1,*)$ by group automorphisms,
 and a nonempty set $\Omega_3$ such that $\Omega$ is of the form
\[(\eas(\Omega_2,\star) \ltimes_\prec \bigl(\eas'(\Omega_1,*)\bigr)\times \eas(\Omega_3),\]
with the products given by
\begin{gather*}
(\alpha_2,\alpha_1,\alpha_3)\rightarrow(\beta_2,\beta_1,\beta_3) =(\alpha_2\star\beta_2,\beta_1\prec\alpha_2,\beta_3),\\
(\alpha_2,\alpha_1,\alpha_3)\triangleright(\beta_2,\beta_1,\beta_3) =\bigl(\alpha_2,\alpha_1*\bigl(\beta_1^{-1}\prec \alpha_2^{-1}\bigr),\alpha_3\bigr).
\end{gather*} \end{Corollary}

\begin{Remark}
The inverse dual CEDS of the CEDS $(\eas(\Omega_1,*) \rtimes_\succ \eas'(\Omega_2,\star))\times \eas(\Omega_3)$ is $(\eas(\Omega_2,\star^{\rm op}) \ltimes_{\succ^{\rm op}} (\eas'(\Omega_1,*))\times \eas(\Omega_3)$.
\end{Remark}


\section{Linear extended associative semigroups}

\subsection{Definitions and examples}

The notions of $\ell$EAS, $\ell$CEDS and dual $\ell$CEDS are introduced in \cite[Definition 1.5]{Foissy48}, as a linear version of Lemma \ref{lem2.6}.

\begin{Definition}\label{defi4.1}
Let $A$ be a vector space and let $\Phi\colon A\otimes A\longrightarrow A\otimes A$ be a linear map.
\begin{enumerate}\itemsep=0pt
\item[(1)] We shall say that $(A,\Phi)$ is a linear extended associative semigroup (briefly, $\ell$EAS) if
\begin{align}
\label{eq30}
(\operatorname{Id} \otimes \Phi)\circ (\Phi\otimes \operatorname{Id})\circ (\operatorname{Id} \otimes \Phi)&=(\Phi \otimes \operatorname{Id})\circ (\operatorname{Id} \otimes \tau)\circ (\Phi\otimes \operatorname{Id}).
\end{align}
\item[(2)] We shall say that $(A,\Phi)$ is a linear commutative extended diassociative semigroup (briefly, $\ell$CEDS) if
\begin{gather}
\tag{\ref{eq30}}(\operatorname{Id} \otimes \Phi)\circ (\Phi\otimes \operatorname{Id})\circ (\operatorname{Id} \otimes \Phi) =(\Phi \otimes \operatorname{Id})\circ (\operatorname{Id} \otimes \tau)\circ (\Phi\otimes \operatorname{Id}),\\
(\operatorname{Id} \otimes \Phi)\circ (\operatorname{Id} \otimes \tau) \circ (\tau \otimes \operatorname{Id})\circ (\Phi \otimes \operatorname{Id})\nonumber\\
\qquad =(\tau \otimes \operatorname{Id})\circ (\Phi \otimes \operatorname{Id})\circ (\operatorname{Id} \otimes \Phi)\circ (\operatorname{Id} \otimes \tau).\label{eq31}
\end{gather}
\item[(3)] We shall say that $(A,\Phi)$ is a linear dual commutative extended diassociative semigroup (briefly, dual $\ell$CEDS) if
\begin{gather}
\tag{\ref{eq30}}(\operatorname{Id} \otimes \Phi)\circ (\Phi\otimes \operatorname{Id})\circ (\operatorname{Id} \otimes \Phi)=(\Phi \otimes \operatorname{Id})\circ (\operatorname{Id} \otimes \tau)\circ (\Phi\otimes \operatorname{Id}),\\
(\Phi \otimes \operatorname{Id})\circ (\tau \otimes \operatorname{Id})\circ (\operatorname{Id} \otimes \tau)\circ (\operatorname{Id} \otimes \Phi)\nonumber\\
\qquad=(\operatorname{Id} \otimes \tau)\circ (\operatorname{Id} \otimes \Phi)\circ (\Phi \otimes \operatorname{Id})\circ (\tau \otimes \operatorname{Id}).\label{eq32}
\end{gather}\end{enumerate}
If $(A,\Phi)$ is an $\ell$EAS (resp.\ an $\ell$CEDS or a dual $\ell$CEDS), we shall say that it is non-degenerate if $\Phi$ is bijective.
\end{Definition}

Note that, by definition, $\ell$CEDS and dual $\ell$CEDS are $\ell$EAS.

\begin{Example}\label{ex4.1} \quad
\begin{enumerate}\itemsep=0pt
\item[(1)] Let $(\Omega,\rightarrow,\triangleright)$ be an EAS (resp.\ a CEDS, a dual CEDS). Let $A=\mathbb{K} \Omega$ be the vector space generated by $\Omega$. We define
\[\Phi\colon \ \begin{cases}
A\otimes A \longrightarrow A\otimes A,\\
a\otimes b \longrightarrow (a\rightarrow b)\otimes (a\triangleright b).
\end{cases}\]
Then $(A,\Phi)$ is an $\ell$EAS (resp.\ an $\ell$CEDS, a dual $\ell$CEDS), which we call the linearization of $(\Omega,\rightarrow,\triangleright)$. It is a non-degenerate $\ell$EAS if and only if $\Omega$ is a non-degenerate EAS.
\item[(2)] Not all the $\ell$EAS are of the form $\mathbb{K}\Omega$. For example, if $A$ is a two-dimensional space with basis $(x,y)$, the maps given by the following matrices in the basis $(x\otimes x,x\otimes y,y\otimes x,y\otimes y)$ are $\ell$EAS
\begin{alignat*}{4}
&M_1 =\begin{pmatrix}0&0&1&0\\0&0&0&0\\0&0&0&0\\0&0&0&0\end{pmatrix},\qquad&&
M_2 =\begin{pmatrix}0&0&0&0\\0&0&a&0\\0&0&0&0\\0&0&0&0\end{pmatrix},\qquad&&
M_3=\begin{pmatrix}1&0&0&0\\0&0&0&0\\0&0&0&0\\0&0&0&0\end{pmatrix},&\\
&M_4=\begin{pmatrix}1&0&0&0\\0&0&1&0\\0&0&0&0\\0&0&0&0\end{pmatrix},\qquad&&
M_5=\begin{pmatrix}1&0&0&0\\0&0&0&0\\0&1&0&0\\0&0&0&0\end{pmatrix},\qquad&&
M_6=\begin{pmatrix}1&0&0&0\\0&0&0&0\\0&0&1&0\\0&0&0&0\end{pmatrix},&\\
&M_7=\begin{pmatrix}1&0&0&0\\0&0&0&0\\0&0&0&0\\0&0&1&0\end{pmatrix},\qquad&&
M_8=\begin{pmatrix}1&0&0&0\\0&0&1&0\\0&0&1&0\\0&0&0&0\end{pmatrix},\qquad&&
M_9=\begin{pmatrix}1&0&0&0\\0&0&0&0\\0&1&0&0\\0&0&1&0\end{pmatrix},&\\
&M_{10}=\begin{pmatrix}1&0&0&0\\0&0&0&0\\0&1&1&0\\0&0&0&0\end{pmatrix},\qquad&&
M_{11}=\begin{pmatrix}1&0&0&0\\0&0&1&0\\0&0&1&0\\0&0&1&0\end{pmatrix},\qquad&&
M_{12}=\begin{pmatrix}1&0&0&0\\0&0&0&0\\0&1&0&0\\0&1&-1&0\end{pmatrix},&\\
&M_{13}=\begin{pmatrix}1&0&0&0\\0&0&0&0\\0&0&0&0\\0&0&0&1\end{pmatrix},\qquad&&
M_{14}=\begin{pmatrix}1&0&0&0\\0&0&0&0\\0&1&0&0\\0&0&0&1\end{pmatrix},\qquad&&
M_{15}=\begin{pmatrix}1&1&0&0\\0&0&0&0\\0&0&0&0\\0&0&1&1\end{pmatrix},&\\
&M_{16}=\begin{pmatrix}1&0&0&0\\0&0&0&0\\0&1&0&0\\0&0&1&1\end{pmatrix},\qquad&&
M_{17}=\begin{pmatrix}1&0&1&0\\0&0&-1&0\\0&1&-1&0\\0&0&2&1\end{pmatrix},\qquad&&
M_{18}=\begin{pmatrix}1&0&0&0\\0&0&1&0\\0&1&0&0\\0&0&0&1\end{pmatrix},&
\end{alignat*}
where $a$ is a scalar. Moreover,
\begin{itemize}\itemsep=0pt
\item The $\ell$CEDS in this list are the $M_i$'s with
\[i\in \{1,2,3,4,5,9,10,13,14,16,17,18\}.\]
\item The dual $\ell$CEDS in this list are the $M_i$'s with
\[i\in \{1,2,3,4,5,7,8,9,11,13,14,15,\allowbreak 16,17,18\}.\]
\end{itemize}
These $\ell$EAS are in fact the EAS of dimension 2 which have a basis of special vectors, see Definition \ref{defi4.3}.
\end{enumerate}\end{Example}

\begin{Notation}
Let $(A,\Phi)$ be an $\ell$EAS. We use Sweedler's-like notation
\[
\Phi(a\otimes b)=\sum a'\rightarrow b' \otimes a''\triangleright b''.
\]
Note that the operations $\rightarrow$ and $\triangleright$ do not necessarily exist, nor the coproducts $a'\otimes a''$ or $b'\otimes b''$.
With this notation, \eqref{eq30} can be rewritten as
\begin{gather}
\sum\sum\sum a' \rightarrow \bigl(b' \rightarrow c'\bigr)'\otimes \bigl(a'' \triangleright \bigl(b' \rightarrow c'\bigr)''\bigr)'\rightarrow \bigl(b'' \triangleright c''\bigr)' \otimes \bigl(a'' \triangleright \bigl(b' \rightarrow c'\bigr)\bigr)''\triangleright \bigl(b'' \triangleright c''\bigr)''\nonumber\\
 \qquad=\sum\sum \bigl(a' \rightarrow b'\bigr)'\rightarrow c' \otimes \bigl(a' \rightarrow b'\bigr)''\triangleright c'' \otimes a'' \triangleright b''.\tag{\ref{eq30}$'$}
\end{gather}
Similarly, \eqref{eq31} and \eqref{eq32} are rewritten as
\begin{gather}
 \sum\sum a'' \triangleright \bigl(c' \rightarrow b'\bigr)'' \otimes a' \rightarrow \bigl(c'' \rightarrow b''\bigr)'\otimes c'' \triangleright b''\nonumber\\
\qquad=\sum\sum a'' \triangleright b'' \otimes c' \rightarrow \bigl(a' \rightarrow b'\bigr)'\otimes c'' \triangleright \bigl(a' \rightarrow b'\bigr)'',\tag{\ref{eq31}$'$}\\
\nonumber\\
 \sum\sum \bigl(b'' \triangleright c''\bigr)'\rightarrow a'\otimes \bigl(b'' \triangleright c''\bigr)''\triangleright a''\otimes b' \rightarrow c'\nonumber\\
\qquad=\sum\sum b' \rightarrow a' \otimes \bigl(b'' \triangleright a''\bigr)''\triangleright c''\otimes \bigl(b'' \triangleright a''\bigr)'\rightarrow c'.\tag{\ref{eq32}$'$}
\end{gather}\end{Notation}

By transposition of \eqref{eq30}, \eqref{eq31} and \eqref{eq32}, we have the following.

\begin{Proposition}\label{prop4.2}
Let $V$ be a finite-dimensional space and
\[
\Phi\colon \ V\otimes V\longrightarrow V\otimes V
\]
 be a linear map. We consider
\[
\Phi^*\colon \ V^*\otimes V^*=(V\otimes V)^*\longrightarrow(V\otimes V)^*=V^*\otimes V^*.
\]
Then $(V,\Phi)$ is an $\ell$EAS $($resp.\ an $\ell$CEDS, a dual $\ell$CEDS$)$ if and only if $(V^*,\Phi^*)$ is an $\ell$EAS $($resp.\ a~dual $\ell$CEDS, an $\ell$CEDS$)$.
\end{Proposition}

\begin{Example}\quad
\begin{enumerate}\itemsep=0pt
\item[(1)] As their matrices are symmetric, the $\ell$EAS $M_3$, $M_6$, $M_{13}$ and $M_{18}$ are self-dual, through the pairing which matrix in the basis $(x,y)$ is
$\bigl(\begin{smallmatrix}
1&0\\0&1
\end{smallmatrix}\bigr)$.
With the same pairing, the dual of $M_4$ is $M_5$ and the dual of $M_8$ is $M_{10}$. The $\ell$EAS $M_2$ and $M_{14}$ are also self-dual, through the pairing which matrix in the basis~$(x,y)$ is
$\bigl(\begin{smallmatrix}
0&1\\1&0
\end{smallmatrix}\bigr)$.
The $\ell$EAS $M_{16}$ and $M_{17}$ are self-dual,\footnote{For $M_{17}$, this holds if the characteristic of the base field is not 2.} through the pairings which matrix in the basis $(x,y)$ are respectively
$\bigl(\begin{smallmatrix}
1&1\\1&0
\end{smallmatrix}\bigr)$, $\bigl(\begin{smallmatrix}
2&2\\2&1
\end{smallmatrix}\bigr)$.
The duals of $M_1$, $M_7$, $M_9$, $M_{11}$, $M_{12}$ and $M_{15}$ are not isomorphic to any $M_i$'s.
\item[(2)] Let $\Omega$ be a finite EAS and $A=\mathbb{K}\Omega$ be the associated $\ell$EAS. The dual $A^*$ is identified with the space $\mathbb{K}^\Omega$ of maps from $\Omega$ to $\mathbb{K}$, with the dual basis $(\delta_\alpha)_{\alpha \in \Omega}$ of the basis~$\Omega$ of~$\mathbb{K}\Omega$.
Then, for any $\alpha,\beta \in \Omega$,
\[
\Phi^*(\delta_\alpha\otimes \delta_\beta)=\sum_{(\gamma,\delta)\in \phi^{-1}(\alpha,\beta)} \delta_\gamma \otimes \delta_\delta.
\]
This is usually not the linearization of an EAS, except if $\Omega$ is non-degenerate: in this case, we recover the linearization of $(\Omega,\curvearrowright,\blacktriangleright)$ of Proposition \ref{prop2.7}.
\end{enumerate}\end{Example}

\subsection{Special vectors, left units and counits}

\begin{Definition}\label{defi4.3}
Let $(A,\Phi)$ be an $\ell$EAS.
\begin{enumerate}\itemsep=0pt
\item[(1)] Let $a\in A$. We shall say that $a$ is a left unit of $(A,\Phi)$ if for any $b\in A$,
\[
\Phi(a\otimes b)=b\otimes a.
\]
\item[(2)] Let $f \in A^*$. We shall say that $f$ is a left counit of $(A,\Phi)$ if
\[
(f \otimes \operatorname{Id})\circ \Phi=\operatorname{Id} \otimes f.
\]
\item[(3)] Let $a\in A$ and $\lambda \in \mathbb{K}$. We shall say that $a$ is a special vector of $(A,\Phi)$ of eigenvalue $\lambda$ if
 \[
 \Phi(a\otimes a)=\lambda a\otimes a.
 \]
\end{enumerate}
\end{Definition}

\begin{Remark}
Let $(A,\Phi)$ be an $\ell$EAS.
\begin{enumerate}\itemsep=0pt
\item[(1)] Any left unit of $(A,\Phi)$ is a special vector of eigenvalue $1$.
\item[(2)] If $A$ is finite-dimensional, its left counits are the left units of $(A^*,\Phi^*)$.
\item[(3)] The set of left units is a subspace of $A$ and the set of left counits a subspace of $A^*$. The set of special vectors of a given eigenvalue is generally not a subspace of $A$.
\end{enumerate}
\end{Remark}

\begin{Lemma}
Let $(A,\Phi)$ be an $\ell$EAS and $a\in A$ be a nonzero special vector of $(A,\Phi)$. Then its eigenvalue $\lambda$ is $0$ or $1$.
\end{Lemma}

\begin{proof}
Let us apply \eqref{eq30} to $a\otimes a\otimes a$. This gives
\[
\lambda^3 a\otimes a\otimes a=\lambda^2 a\otimes a\otimes a.
\] As $a\neq 0$, $\lambda=0$ or $1$.\end{proof}

\begin{Example}\quad
\begin{enumerate}\itemsep=0pt
\item[(1)] Let us give special vectors, left units and left counits for the thirteen $\ell$EAS associated to the thirteen EAS of cardinality 2. In each case, we give a basis of the spaces of left units and left counits; $\lambda$, $\mu$ and $\nu$ are scalars.
The dual basis of the basis $(X,Y)$ of $\mathbb{K}\Omega$ is denoted by $(X^*,Y^*)$.
$$\renewcommand{\arraystretch}{1.2} \begin{array}{|c||c|c|c|c|}
\hline \text{Case}&\text{Special vectors}&\text{Special vectors}&\text{Left units}&\text{Left conits}\\
&\text{of eigenvalue 1}&\text{of eigenvalue 0}&&\\
\hline \hline \mathbf{A1}&\lambda X&\nu(X-Y)&\varnothing&\varnothing\\
\hline \mathbf{A2}&\lambda X&\nu(X-Y)&\varnothing&(X^*+Y^*)\\
\hline \mathbf{C1}&\lambda X&0&\varnothing&\varnothing\\
\hline \mathbf{C3}&\lambda X,\, \mu Y&0&(Y)&(X^*+Y^*)\\
\hline \mathbf{C5}&\mu Y&0&(Y)&\varnothing\\
\hline \mathbf{C6}&\lambda X&0&\varnothing&\varnothing\\
\hline \mathbf{E'1}\text{--}\mathbf{E'2}&\lambda X&\nu(X-Y)&\varnothing&\varnothing\\
\hline \mathbf{E'3}&\lambda X,\, \mu Y&\nu(X-Y)&\varnothing&(X^*+Y^*)\\
\hline \mathbf{F1}&\lambda X&\nu(X-Y)&(X)&\varnothing\\
\hline \mathbf{F3}&\lambda X+\mu Y&0&(X,Y)&(X^*,Y^*)\\
\hline \mathbf{F4}&\lambda X,\, \nu(X+Y)&0&(X+Y)&(X^*)\\
\hline \mathbf{H1}&\lambda X&0&(X)&\varnothing\\
\hline \mathbf{H2}&\lambda X,\, \nu(X+Y)&0&(X)&(X^*+Y^*)\\
\hline \end{array}$$
Some of them have a basis of special vectors: let us determine their matrices in such a~basis. We recover in this way some matrices of Example~\ref{ex4.1}:
\begin{itemize}\itemsep=0pt
\item For $\mathbf{A1}$, in the basis $(X,X-Y)$, we obtain $M_3$.
\item For $\mathbf{A2}$, in the basis $(X,X-Y)$, and for $\mathbf{F1}$, in the basis $(X-Y,X)$, we obtain $M_4$.
\item For $\mathbf{C3}$, in the basis $(Y,X)$, we obtain $M_{16}$.
\item For $\mathbf{E'1}\text{--}\mathbf{E'2}$, in the basis $(X,X-Y)$, we obtain $M_6$.
\item For $\mathbf{E'3}$, in the basis $(X,Y-X)$, we obtain $M_{11}$.
\item For $\mathbf{F3}$, in the basis $(X,Y)$, we obtain $M_{18}$.
\item For $\mathbf{F4}$ and $\mathbf{H2}$, in the basis $(X-Y,X)$, we obtain $M_{17}$.
\end{itemize}
Hence, the $\ell$EAS associated to $\mathbf{A2}$ and $\mathbf{F1}$ are isomorphic, whereas the EAS $\mathbf{A2}$ and $\mathbf{F1}$ are not. As similar situation holds for $\mathbf{F4}$ and $\mathbf{H2}$.

\item[(2)] It is possible to show that any 2-dimensional $\ell$EAS with a basis of special vectors is isomorphic to one of the eighteen cases of Example \ref{ex4.1}.
For all of them, let us give special vectors, left units and left counits for the eighteen cases of Example~\ref{ex4.1}.
In each case, we give a basis of the spaces of left units and left counits; $\lambda$, $\mu$ and $\nu$ are scalars. For~$M_2$, we~assume that the parameter $a$ is nonzero (otherwise, $\Phi=0$).
$$\renewcommand{\arraystretch}{1.2} \begin{array}{|c||c|c|c|c|}
\hline \text{Case}&\text{Special vectors}&\text{Special vectors}&\text{Left units}&\text{Left counits}\\
&\text{of eigenvalue 1}&\text{of eigenvalue 0}&&\\
\hline\hline M_1&0&\lambda x, \mu y&\varnothing&\varnothing\\
\hline M_2&0&\lambda x, \mu y&\varnothing&\varnothing\\
\hline M_3&\lambda x&\mu y&\varnothing&\varnothing\\
\hline M_4&\lambda x&\mu y&\varnothing& (x^*)\\
\hline M_5&\lambda x&\mu y&(x)&\varnothing\\
\hline M_6&\lambda x&\mu y&\varnothing&\varnothing\\
\hline M_7&\lambda x&\mu y&\varnothing&\varnothing\\
\hline M_8&\lambda x&\mu y&\varnothing& (x^*)\\
\hline M_9&\lambda x&\mu y&(x)&\varnothing\\
\hline M_{10}&\lambda x&\mu y&(x)&\varnothing\\
\hline M_{11}&\lambda x,\nu(x+y)&\mu y&\varnothing& (x^*)\\
\hline M_{12}&\lambda x&\mu y&\varnothing&\varnothing\\
\hline M_{13}&\lambda x, \mu y&0&\varnothing&\varnothing\\
\hline M_{14}&\lambda x, \mu y&0&(x)&(y^*)\\
\hline M_{15}&\lambda x, \mu y&\nu(x-y)&\varnothing&(x^*+y^*)\\
\hline M_{16}&\lambda x, \mu y&0&(x)&(x^*+y^*)\\
\hline M_{17}&\lambda x, \mu y&0&(x)&(x^*+y^*)\\
\hline M_{18}&\lambda x+\mu y&0&(x,y)&(x^*,y^*)\\
\hline \end{array}$$
Among them, $M_{11}$ has three lines of special vectors. In the basis $(x+y,x)$ its matrix is~$M_{15}$, so $M_{11}$ and $M_{15}$ are isomorphic.
\end{enumerate}\end{Example}


\subsection{Left units and counits of finite non-degenerate CEDS}

\begin{Proposition}\label{prop4.5}
Let $(\Omega_1,*)$ be an abelian finite group, $(\Omega_2,\star)$ be a finite group, and $\Omega_3$ be a~finite set. We denote by $e_1$ and $e_2$ the units of $\Omega_1$ and $\Omega_2$.
\begin{enumerate}\itemsep=0pt
\item[$(1)$] Let $(A,\Phi)$ be the linearization of the CEDS
$(\eas(\Omega_1,*) \rtimes_\succ \eas'(\Omega_2,\star))\times \eas(\Omega_3)$
of Theorem {\rm\ref{theo3.14}}.
\begin{enumerate}\itemsep=0pt
\item[$(a)$] The special vectors of eigenvalue $1$ of $(A,\Phi)$ are the vectors of the form
\[a=\sum_{(\alpha_1,\alpha_2,\alpha_3)\in H_1\times H_2\times \Omega_3}g(\alpha_3)(\alpha_1,\alpha_2,\alpha_3),\]
where $H_1$ is a subgroup of $\Omega_1$, $H_2$ is a subgroup of $\Omega_2$, such that $H_2\succ H_1\subseteq H_1$, and~${g\colon \Omega_3\longrightarrow \mathbb{K}}$ is a map.
\item[$(b)$] The left units of $(A,\Phi)$ are the vectors of the form
\[a=\sum_{(\alpha_2,\alpha_3)\in \Omega_2\times \Omega_3}g(\alpha_3)(e_1,\alpha_2,\alpha_3),\]
where $g\colon \Omega_3\longrightarrow \mathbb{K}$ is a map.
\item[$(c)$] The left counits of $(A,\Phi)$ are the linear forms $f$ such that for any $(\alpha_1,\alpha_2,\alpha_3)\in \Omega$,
$f(\alpha_1,\alpha_2,\alpha_3)=\delta_{\alpha_2,e_2}g(\alpha_3)$,
where $g\colon \Omega_3\longrightarrow \mathbb{K}$ is a map.
\end{enumerate}
\item[$(2)$] Let $(A,\Phi)$ be the linearization of the CEDS
$(\eas(\Omega_2,\star) \ltimes_\prec \eas'(\Omega_1,*))\times \eas(\Omega_3)$
of Corollary {\rm\ref{cor3.15}}.
\begin{enumerate}\itemsep=0pt
\item[$(a)$] The special vectors of eigenvalue 1 of $(A,\Phi)$ are the vectors of the form
\[a=\sum_{(\alpha_1,\alpha_2,\alpha_3)\in H_1\times H_2\times \Omega_3}g(\alpha_3)(\alpha_2,\alpha_1,\alpha_3),\]
where $H_1$ is a subgroup of $\Omega_1$, $H_2$ is a subgroup of $\Omega_2$, such that $H_1\prec H_2\subseteq H_1$, and~${g\colon \Omega_3\longrightarrow \mathbb{K}}$ is a map.
\item[$(b)$] The left units of $(A,\Phi)$ are the vectors of the form
\[a=\sum_{(\alpha_1,\alpha_3)\in \Omega_1\times \Omega_3}g(\alpha_3)(e_2,\alpha_1,\alpha_3),\]
where $g\colon \Omega_3\longrightarrow \mathbb{K}$ is a map.
\item[$(c)$] The left counits of $(A,\Phi)$ are the linear forms $f$ such that for any $(\alpha_1,\alpha_2,\alpha_3)\in \Omega$,
$f(\alpha_1,\alpha_2,\alpha_3)=\delta_{\alpha_1,e_1}g(\alpha_3)$,
where $g\colon \Omega_3\longrightarrow \mathbb{K}$ is a map.
\end{enumerate} \end{enumerate}\end{Proposition}

\begin{proof}
(1)(a) Let $a$ be a nonzero vector of $A$, which we write as
\[a=\sum_{(\alpha_1,\alpha_2,\alpha_3)\in \Omega}a_{(\alpha_1,\alpha_2,\alpha_3)}(\alpha_1,\alpha_2,\alpha_3).\]
Then $a$ is a special vector of eigenvalue~1 if and only if for any $(\alpha_1,\alpha_2,\alpha_3)$, $(\beta_1,\beta_2,\beta_3)\in \Omega$,
\begin{align}
\label{eq33}a_{(\alpha_1,\alpha_2,\alpha_3)}a_{(\beta_1,\beta_2,\beta_3)}&=a_{(\alpha_2^{-1}\succ \beta_1,\beta_2\star \alpha_2,\beta_3)}a_{((\alpha_2^{-1}\succ \beta_1^{-1})*\alpha_1,\alpha_2,\alpha_3)}.
\end{align}
We put
\begin{gather*}
 \forall (\alpha_2,\alpha_3)\in \Omega_2\times \Omega_3,\qquad H_1(\alpha_2,\alpha_3)=\{\alpha_1\in \Omega_1, a_{(\alpha_1,\alpha_2,\alpha_3)}\neq 0\},\\
 \forall \alpha_3\in \Omega_3,\qquad H_2(\alpha_3)=\{\alpha_2\in \Omega_2, H_1(\alpha_2,\alpha_3)\neq \varnothing\},\qquad
 H_3=\{\alpha_3\in \Omega_3, H_2(\alpha_3)\neq \varnothing\}.
\end{gather*}
We shall also consider the map
\[g\colon \ \begin{cases}
\Omega_3 \longrightarrow \mathbb{K},\\
\alpha_3 \longrightarrow g(\alpha_3)=a_{e_1,e_2,\alpha_3}.
\end{cases}\]

Let us first prove that if $\alpha_3\in H_3$, then $H_2(\alpha_3)$ is a subgroup of $\Omega_2$. Let $\alpha_2,\beta_2\in H_2(\alpha_3)$ (which is nonempty as $\alpha_3\in H_3$).
Let $\alpha_1,\beta_1\in \Omega_1$, such that $a_{(\alpha_1,\alpha_2,\alpha_3)}\neq 0$ and $a_{(\beta_1,\beta_2,\alpha_3)}\neq 0$. By \eqref{eq33},
\smash{$a_{(\alpha_2^{-1}\succ \beta_1,\beta_2\star \alpha_2,\alpha_3)}\neq 0$},
so $\beta_2\star \alpha_2\in H_2(\alpha_3)$. As $\Omega_2$ is finite, $H_2(\alpha_3)$ is a subgroup of~$\Omega_2$.

Let us prove that if $\alpha_3\in H_3$, then $H_1(e_2,\alpha_3)$ is a subgroup of $\Omega_1$ and, moreover, for any $\alpha_1\in H_1(e_2,\alpha_3)$, $a_{(\alpha_1,e_2,\alpha_3)}=g(\alpha_3)$.
As $H_2(\alpha_3)$ is a subgroup of $\Omega_2$, it contains $e_2$, so $H_1(e_2,\alpha_3)\neq \varnothing$. Let $\alpha_1,\beta_1\in H_1(e_2,\alpha_3)$. By \eqref{eq33},
\[a_{(\alpha_1,e_2,\alpha_3)}a_{(\beta_1,e_2,\alpha_3)}=a_{(\beta_1,e_2,\alpha_3)}a_{(\beta_1^{-1}*\alpha_1,e_2,\alpha_3)}\neq 0.\]
Hence, $\beta_1^{-1}*\alpha_1\in H_1(e_2,\alpha_3)$. Taking $\alpha_1=\beta_1$, we obtain{\samepage
\[a_{(\alpha_1,e_2,\alpha_3)}a_{(\alpha_1,e_2,\alpha_3)}=a_{(\alpha_1,e_2,\alpha_3)}a_{(e_1,e_2,\alpha_3)}\neq 0,\]
so $a_{(\alpha_1,e_2,\alpha_3)}=a_{(e_1,e_2,\alpha_3)}=g(\alpha_3)$.}

Let us prove that if $\alpha_3,\beta_3\in H_3$ and $\beta_2\in H_2(\beta_3)$, then $H_1(\beta_2,\beta_3)\subseteq H_1(e_2,\alpha_3)$. Let $\beta_1\in H_1(\beta_2,\beta_3)$.
Then $a_{(\beta_1,\beta_2,\beta_3)}\neq 0$. As $H_1(e_1,\beta_2)$ is a subgroup of $\Omega_1$, it contains $e_1$, so~${a_{(e_1,e_2,\alpha_3)}\neq 0}$. By \eqref{eq33},
\[a_{(e_1,e_2,\alpha_3)}a_{(\beta_1,\beta_2,\beta_3)}=g(\alpha_3)a_{(\beta_1,\beta_2,\beta_3)}=a_{(\beta_1,\beta_2,\beta_3)}a_{(\beta_1^{-1},e_2,\alpha_3)}\neq 0,\]
so $\beta_1^{-1}\in H_1(e_2,\alpha_3)$. As this is a subgroup of $\Omega_1$, $\beta_1\in H_1(e_2,\alpha_3)$.
As a consequence, for $\beta_2=e_2$, we obtain by symmetry that for any $\alpha_3,\beta_3\in H_3$, $H_1(e_2,\alpha_3)=H_1(e_2,\beta_3)$.
Therefore, there exists a subgroup $H_1$ of $\Omega_1$ such that for any $\alpha_3\in \Omega_3$, $H_1(e_2,\alpha_3)=H_1$.

Let us prove that for any $\alpha_3\in H_3$, $H_2(\alpha_3)\succ H_1\subseteq H_1$. Let $\beta_1\in H_1=H_1(e_2,\alpha_3)$, then $a_{(\beta_1,e_2,\alpha_3)}\neq 0$. Let $\alpha'_2\in H_2(\alpha_3)$.
We put $\alpha_2=\alpha_2'^{-1}\in H_2(\alpha_3)$. There exists $\alpha_1\in H_1(\alpha_2,\alpha_3)$, such that $a_{(\alpha_1,\alpha_2,\alpha_3)}\neq 0$. By \eqref{eq33},
 \[a_{(\alpha_1,\alpha_2,\alpha_3)}a_{(\beta_1,e_2,\alpha_3)}=a_{(\beta_1,\alpha_2,\alpha_3)}a_{((\alpha_2^{-1}\succ \beta_1^{-1})*\alpha_1,e_2,\alpha_3)}\neq 0,\]
so $\bigl(\alpha_2^{-1}\succ \beta_1^{-1}\bigr)*\alpha_1\in H_1(e_2,\alpha_3)=H_1$. Moreover, as $H_1(\alpha_2,\alpha_2)\subseteq H_1$,
\[\alpha_2^{-1}\succ \beta_1^{-1}=\alpha_2'\succ \beta_1^{-1}\in H_1.\]
Its inverse $ \alpha'_2\succ \beta_1$ is also an element of $H_1$, so $H_2(\alpha_3)\succ H_1\subseteq H_1$.

Let us prove that for any $\alpha_3\in H_3$, for any $\alpha_2\in H_2(\alpha_3)$, $H_1(\alpha_2,\alpha_3)=H_1$. Let $\alpha_1\in H_1=H_1(e_2,\alpha_3)$ and $\beta_1\in H_1(\alpha_2,\alpha_3)$. By \eqref{eq33},
\[a_{(\alpha_1,e_2,\alpha_3)}a_{(\beta_1,\alpha_2,\alpha_3)}=a_{(\beta_1,\alpha_2,\alpha_3)}a_{(\beta_1^{-1}*\alpha_1,\alpha_2,\alpha_3)}\neq 0,\]
so $\beta_1^{-1}*\alpha_1\in H_1(\alpha_2,\alpha_3)$. We obtain an injective map
\[\begin{cases}
H_1 \longrightarrow H_1(\alpha_2,\alpha_3),\\
\alpha_1 \longrightarrow \beta_1^{-1}*\alpha_1.
\end{cases}\]
Hence, $|H_1|\leqslant |H_1(\alpha_2,\alpha_3)|$. We already proved that $H_1(\alpha_2,\alpha_3)\subseteq H_1$, so $H_1=H_1(\alpha_2,\alpha_3)$.

We now prove that there exists a subgroup $H_2$ of $\Omega_2$ such that for any $\alpha_3\in H_3$, $H_2(\alpha_3)=H_2$, and that, moreover, for any $\alpha_2\in H_2$, for any $\alpha_3\in H_3$, $a_{(e_1,\alpha_2,\alpha_3)}=g(\alpha_3)$.
Let $\alpha_3,\beta_3\in H_3$. Let $\alpha_2\in H_2(\alpha_3)$. As $e_1\in H_1(\alpha_2,\alpha_3)=H_1$, $a_{(e_1,\alpha_2,\alpha_3)}\neq 0$. As $e_2\in H_2(\alpha_3)$, $a_{(e_1,e_2,\beta_3)}\neq 0$. By~\eqref{eq33},
\[a_{(e_1,\alpha_2,\alpha_3)}a_{(e_1,e_2,\beta_3)}=a_{(e_1,\alpha_2,\beta_3)}a_{(e_1,\alpha_2\alpha_3)}\neq 0,\]
so $\alpha_2\in H_2(\beta_3)$. We proved that $H_2(\alpha_3)\subseteq H_2(\beta_3)$: by symmetry, $H_2(\alpha_3)=H_2(\beta_3)$, which prove the existence of $H_2$.
Moreover, as $a_{(e_1,\alpha_2,\alpha_3)}\neq 0$ and $a_{(e_1,e_2,\beta_3)}=g(\beta_3)$, we obtain that $a_{(e_1,\alpha_2,\beta_3)}=g(\beta_3)$.

We proved that for any $(\alpha_1,\alpha_2,\alpha_3)\in \Omega$,
\[a_{(\alpha_1,\alpha_2,\alpha_3)}\neq 0\Longleftrightarrow (\alpha_1,\alpha_2,\alpha_3)\in H_1\times H_2\times H_3.\]
Let $\alpha_3,\beta_3\in H_3$, $\beta_1\in H_1$, $\alpha_2\in H_2$. We put $\alpha_1=\alpha_2^{-1}\succ \beta_1$. Then, by \eqref{eq33},
\[g(\alpha_3)g(\beta_3)=a_{(\beta_1,\alpha_2,\alpha_3)}a_{(e_1,\alpha_2,\beta_3)}=a_{(\beta_1,\alpha_2,\alpha_3)}g(\beta_3),\]
so $a_{(\beta_1,\alpha_2,\alpha_3)}=g(\alpha_3)$. We proved that $a$ has the announced form.

Conversely, if $a$ is of the announced form,
\begin{align*}
\Phi(a\otimes a)={}&\sum_{(\alpha_1,\alpha_2,\alpha_3),(\beta_1,\beta_2,\beta_3)\in H_1\times H_2\times \Omega_3}g(\alpha_3)g(\beta_3)(\alpha_1*\beta_1,\beta_2,\beta_3)\\
&\otimes \bigl(\beta_2\succ \alpha_2,\alpha_2*\beta_2^{-1},\alpha_3\bigr)\\
={}&\sum_{(\alpha_1,\alpha_2,\alpha_3),(\beta_1,\beta_2,\beta_3)\in H_1\times H_2\times \Omega_3}(\beta_1,\beta_2,\beta_3)\otimes (\alpha_1,\alpha_2,\alpha_3)
=a\otimes a.
\end{align*}

(2)(a) If $(A,\Phi)$ is the linearization of $(\eas(\Omega_2,\star) \ltimes_\prec \eas'(\Omega_1,*))\times \eas(\Omega_3)$, then $\bigl(A,\Phi^{-1}\bigr)$ is the linearization of $(\eas(\Omega_1,*) \rtimes_{\prec^{\rm op}} \eas'(\Omega_2,\star^{\rm op}))\times \eas(\Omega_3)$.
The result then comes from the observation that the special vectors of $(A,\Phi)$ and $\bigl(A,\Phi^{-1}\bigr)$ are the same.

(1)(b) Let $a$ be a left unit of $A$. Then it is a special vector, which we write as
\[a=\sum_{(\alpha_1,\alpha_1,\alpha_3)\in H_1\times H_2\times \Omega_3}g(\alpha_3)(\alpha_1,\alpha_1,\alpha_3).\]
For any $b=(\beta_1,\beta_2,\beta_3)\in \Omega$,
\begin{align*}
\Phi(a\otimes b)&=\sum_{(\alpha_1,\alpha_1,\alpha_3)\in H_1\times H_2\times \Omega_3}g(\alpha_3)(\alpha_1*\beta_1,\beta_2,\beta_3)\otimes \bigl(\beta_2\succ \alpha_1,\alpha_2*\beta_2^{-1},\alpha_3\bigr)\\
&=b\otimes a=\sum_{(\alpha_1,\alpha_1,\alpha_3)\in H_1\times H_2\times \Omega_3}g(\alpha_3)(\beta_1,\beta_2,\beta_3)\otimes (\alpha_1,\alpha_1,\alpha_3).
\end{align*}
Taking $\beta_1=e_1$, we obtain that for any $\alpha_1\in H_1$, $\alpha_1=e_1$, so $H_1=\{e_1\}$. Moreover, for any $\beta_2\in \Omega_2$,
\[\sum_{(\alpha_2,\alpha_3)\in H_2\times \Omega_3}g(\alpha_3)\bigl(\alpha_2\star \beta_2^{-1},\alpha_3\bigr)=\sum_{(\alpha_2,\alpha_3)\in H_2\times \Omega_3}g(\alpha_3)(\alpha_2,\alpha_3),\]
so for any $\alpha_2\in H_2$, $\alpha_2\star \beta_2^{-1}\in H_2$. In particular, for $\alpha_2=e_2$, $\beta_2^{-1}\in H_2$ and finally $\beta_2\in H_2$: $H_2=\Omega_2$. The converse application is immediate.

(2)(b) Similar proof.

(1)(c) and (2)(c) The left counits of $(A,\Phi)$ are the left units of $(A^*,\Phi^*)$, which is isomorphic to $\bigl(A,\Phi^{-1}\bigr)$.The result comes from(2)(b). and (1)(b). \end{proof}


\section[From bialgebras to lEAS]{From bialgebras to $\boldsymbol{\ell}$EAS}

We refer to \cite{Abe1980,Cartier2021,Sweedler1969} for classical results and notations on bialgebras and Hopf algebras.

\subsection[A functor from bialgebras to lEAS]{A functor from bialgebras to $\boldsymbol{\ell}$EAS}

\begin{Proposition} \label{prop5.1}
Let $(A,m,\Delta)$ be a bialgebra, not necessarily unitary nor counitary. For any $a,b\in A$, we define $\Phi\colon A\otimes A\longrightarrow A\otimes A$ by
\begin{align*}
\Phi(a\otimes b) =(m\otimes \operatorname{Id}_A)\circ (\operatorname{Id}_A\otimes \tau)\circ (\Delta\otimes \operatorname{Id}_A)(a\otimes b)=\sum a^{(1)}b\otimes a^{(2)},
\end{align*}
with Sweedler's notation $\Delta(a)=\sum a^{(1)}\otimes a^{(2)}$. Then $(A,\Phi)$ is an $\ell$EAS, denoted by $\leas(A,\allowbreak m,\Delta)$.
\end{Proposition}

\begin{proof}
For any $a,b,c\in A$,
\begin{align*}
(\operatorname{Id} \otimes \Phi)\circ(\Phi\otimes \operatorname{Id})\circ (\operatorname{Id} \otimes \Phi)(a\otimes b\otimes c)&=(\Phi\otimes \operatorname{Id})\circ (\operatorname{Id} \otimes \tau)\circ (\Phi \otimes \operatorname{Id})(a\otimes b\otimes c)\\
&=\sum \sum a^{(1)}b^{(1)}c\otimes a^{(2)}b^{(2)}\otimes a^{(3)}. \tag*{\qed}
\end{align*} \renewcommand{\qed}{}\end{proof}

\begin{Example}\quad\samepage
\begin{enumerate}\itemsep=0pt
\item[(1)] Let $(\Omega,\star)$ be a semigroup. We take $A=\mathbb{K}\Omega$, with its usual bialgebra structure: the product $m$ obtained by linearization of $\star$ and the coproduct $\Delta$ defined by
$\forall \alpha \in \Omega$, $\Delta(\alpha) =\alpha \otimes \alpha$.
Then $(A,m,\Delta)$ is a counitary bialgebra, unitary if and only if $\Omega$ is a monoid. In~$\leas(A,m,\Delta)$, for any $\alpha,\beta \in \Omega$,
$\Phi(\alpha \otimes \beta)=\alpha\star \beta \otimes \alpha$.
We recover the linearization of $\eas(\Omega,\star)$.
\item[(2)] Let $A$ be a vector space, $1_A\in A$ and $\varepsilon\in A^*$ such that $\varepsilon(1_A)=1$. We define a product and a coproduct on $A$ by
$
\forall a,b\in A$, $ a\cdot b =\varepsilon(a) b$, $
\forall a\in A$, $ \Delta(a) =1_A\otimes a$.
Then $(A,m,\Delta)$ is a bialgebra, with a left unit $1_A$ and a left counit $\varepsilon$. It is unitary if and only if $A$ is one-dimensional; it is counitary if, and only if, $A$ is one-dimensional. In $\leas(A,m,\Delta)$, for any $a,b\in A$, $\Phi(a\otimes b)=b\otimes a$.
\end{enumerate}\end{Example}

\begin{Proposition}\label{prop5.2}
Let $(A,m,\Delta)$ be a bialgebra, not necessarily unitary nor counitary.
\begin{enumerate}\itemsep=0pt
\item[$(1)$] Let us consider the following conditions:
\begin{enumerate}\itemsep=0pt
\item[$(a)$] $\leas(A,m,\Delta)$ is an $\ell$CEDS.
\item[$(b)$] For any $a,b,c\in A$, $\sum \sum a^{(1)}b^{(1)}c\otimes a^{(2)}\otimes b^{(2)}=\sum \sum b^{(1)}a^{(1)}c\otimes a^{(2)}\otimes b^{(2)}$.
\item[$(c)$] For any $a,b,c\in A$, $abc=bac$.
\item[$(d)$] $m$ is commutative.
\end{enumerate}
Then $(d)\Longrightarrow (c)\Longrightarrow (b)\Longleftrightarrow (a)$. If $(A,\Delta)$ has a right counit, then $(c)\Longleftrightarrow (a)$. If $(A,m,\Delta)$ has a right counit and a right unit, then $(d)\Longleftrightarrow (a)$.
\item[$(2)$] Let us consider the following conditions:
\begin{enumerate}\itemsep=0pt
\item[$(a)$] $\leas(A,m,\Delta)$ is a dual $\ell$CEDS.
\item[$(b)$] For any $a,b,c\in A$,
$\sum a^{(1)}b\otimes a^{(2)}c\otimes a^{(3)}=\sum a^{(2)}b\otimes a^{(1)}c \otimes a^{(3)}$.
\item[$(c)$] $(\Delta \otimes \operatorname{Id})\circ \Delta=(\tau \otimes \operatorname{Id})\circ (\Delta \otimes \operatorname{Id})\circ \Delta$.
\item[$(d)$] $\Delta$ is cocommutative.
\end{enumerate}
Then $(d)\Longrightarrow (c)\Longrightarrow (b)\Longleftrightarrow (a)$. If $(A,m)$ has a right unit, then $(c)\Longleftrightarrow (a)$. If $(A,m,\Delta)$ has a right counit and a right unit, then $(d)\Longleftrightarrow (a)$.
\end{enumerate}
\end{Proposition}

\begin{proof}
(1) Obviously, $(d)\Longrightarrow (c)\Longrightarrow (b)$. Let $a,b,c\in A$. Then
\begin{align*}
(\operatorname{Id} \otimes \Phi)\circ (\operatorname{Id} \otimes \tau)\circ (\tau \otimes \operatorname{Id})\circ (\Phi\otimes \operatorname{Id})(b\otimes c\otimes a)&=\sum \sum b^{(2)}\otimes a^{(1)}b^{(1)}c\otimes a^{(2)},\\
(\tau \otimes \operatorname{Id})\circ (\Phi\otimes \operatorname{Id})\circ (\operatorname{Id} \otimes \Phi)\circ (\operatorname{Id} \otimes \tau)(b\otimes c\otimes a)&=\sum \sum b^{(2)}\otimes b^{(1)}a^{(1)}c\otimes a^{(2)},
\end{align*}
so $(a)\Longleftrightarrow (b)$. If $(b)$ is satisfied and if $(A,\Delta)$ has a right counit $\varepsilon$, applying $(\operatorname{Id} \otimes \varepsilon \otimes \varepsilon)$ to $(b)$, we obtain $(c)$.
If $(c)$ is satisfied and $(A,m)$ has a right unit $1_A$, taking $c=1_A$ in $(c)$, we obtain~$(d)$.

(2) Obviously, $(d)\Longrightarrow (c)\Longrightarrow (b)$. Let $a,b,c\in A$. Then
\begin{align*}
(\Phi\otimes \operatorname{Id})\circ (\tau \otimes \operatorname{Id})\circ (\operatorname{Id} \otimes \tau)\circ (\operatorname{Id} \otimes \Phi)(b\otimes a\otimes c)&=\sum a^{(2)}b\otimes a^{(3)}\otimes a^{(1)}c,\\
(\operatorname{Id} \otimes \tau)\circ (\operatorname{Id} \otimes \Phi)\circ (\Phi\otimes \operatorname{Id})\circ (\tau \otimes \operatorname{Id})(b\otimes a\otimes c)&=\sum a^{(1)}b\otimes a^{(3)}\otimes a^{(2)}c,
\end{align*}
so $(a)\Longleftrightarrow (b)$. If $(b)$ is satisfied and if $(A,m)$ has a right unit $1_A$, taking $b=c=1_A$ in $(b)$, we~obtain~$(c)$. If $(c)$ is satisfied
and $(A,\Delta)$ has a~right counit $\varepsilon$, applying $(\operatorname{Id} \otimes \operatorname{Id} \otimes \varepsilon)$ to $(c)$, we~obtain~$(d)$. \end{proof}

\begin{Proposition}
Let $(A,m,\Delta)$ be a finite-dimensional bialgebra, not necessarily unitary nor counitary. Then $\leas(A,m,\Delta)^*=\leas(A^*,\Delta^*,m^*)$.
\end{Proposition}

\begin{proof}
Let $f,g\in A^*$. For any $a,b\in A$,
\begin{align*}
\Phi^*&=((m\otimes \operatorname{Id}_A)\circ (\operatorname{Id}_A\otimes \tau)\circ (\Delta\otimes \operatorname{Id}_A) )^*=(\Delta^*\otimes \operatorname{Id}_{A^*})\circ (\operatorname{Id}_{A^*}\otimes \tau)\circ (m^*\otimes \operatorname{Id}_{A^*}).
\end{align*}
Therefore, $\leas(A,m,\Delta)^*=\leas(A^*,\Delta^*,m^*)$. \end{proof}

\begin{Proposition}\label{prop5.4}
Let $(A,m,\Delta)$ be a bialgebra.
\begin{enumerate}\itemsep=0pt
\item[$(1)$] We assume that $(A,m)$ has a right unit $1_A$.
\begin{itemize}\itemsep=0pt
\item If $1_A$ is not a unit of $(A,m)$, the unique left unit of $\leas(A,m,\Delta)$ is $0$. If $1_A$ is a unit of $(A,m)$, then the left units of $\leas(A,m,\Delta)$ are the elements $a\in A$ such that $\Delta(a)=1_A\otimes a$.
\end{itemize}
\item[$(2)$] We assume that $(A,\Delta)$ has a right counit $\varepsilon_A$.
\begin{itemize}\itemsep=0pt
\item If $\varepsilon_A$ is not a unit of $(A,\Delta)$, the unique left counit of $\leas(A,m,\Delta)$ is $0$. If $\varepsilon_A$ is a counit of $(A,\Delta)$, then the left counits of $\leas(A,m,\Delta)$ are the elements $\lambda \in A^*$
such that $\lambda \circ m=\varepsilon\otimes \lambda$.
\end{itemize}\end{enumerate}\end{Proposition}

\begin{proof}
(1) Let us assume that $\leas(A,m,\Delta)$ has a nonzero left unit $a$. Let us choose $\lambda \in A^*$ such that $\lambda(a)=1$. For any $b\in A$,
\begin{align*}
(\operatorname{Id} \otimes \lambda)\circ \Phi(a\otimes b)&=\underbrace{\left(\sum a^{(1)}\lambda\bigl(a^{(2)}\bigr)\right)}_{=a'} b=(\operatorname{Id} \otimes \lambda)(b\otimes a)=b\lambda(a)=b,
\end{align*}
so $a'$ is a left unit of $(A,m)$. Then $a'1_A=a'=1_A$, so $a'=1_A$ is a unit. Moreover, for $b=1_A$,
\[\Phi(a\otimes 1_A)=\sum a^{(1)}1_A\otimes a^{(2)}=\Delta(a)=1_A\otimes a.\]
Conversely, if $1_A$ is a unit of $(A,m)$ and $\Delta(a)=1_A\otimes a$, then $a$ is clearly a left unit of $\leas(A,m,\Delta)$.

(2) Let us assume that $\leas(A,m,\Delta)$ has a nonzero left counit $\lambda$. Let us choose $b\in A$ such that $\lambda(b)=1$. For any $a\in A$,
\begin{align*}
(\lambda \otimes \operatorname{Id})\circ \Phi(a\otimes b)&=\sum \lambda\bigl(a^{(1)}b\bigr)a^{(2)}=a\lambda(b)=a.
\end{align*}
If we define $\lambda'\colon A\longrightarrow \mathbb{K}$ by $\lambda'(a)=\lambda(ab)$, then $\lambda'$ is a left counit of $(A,\Delta)$. As $\varepsilon_A$ is a right counit of $(A,\Delta)$,
$(\lambda'\otimes \varepsilon_A)\circ \Delta=\lambda'=\varepsilon_A$,
so $\lambda'=\varepsilon_A$ is a counit of $(A,\Delta)$. Moreover, for any~${a,b\in A}$,
\begin{align*}
(\lambda \otimes \varepsilon_A)\circ \Phi(a\otimes b)&=\varepsilon_A(a)\lambda(b)=\sum \lambda\bigl(a^{(1)}b\bigr)\varepsilon_A\bigl(a^{(2)}\bigr)=\lambda(ab),
\end{align*}
so $\lambda \circ m=\varepsilon_A\otimes \lambda$. Conversely, if $\varepsilon_A$ is a counit of $(A,\Delta)$ and $\lambda \circ m=\varepsilon_A\otimes \lambda$, then for any $a,b\in A$,
\begin{align*}
(\lambda \otimes \operatorname{Id})\circ \Phi(a\otimes b)&=\sum \lambda\bigl(a^{(2)}b\bigr)a^{(2)}=\sum \varepsilon_A\bigl(a^{(1)}\bigr)\lambda(b)a=\lambda(b)a,
\end{align*}
so $\lambda$ is a left counit of $\leas(A,m,\Delta)$. \end{proof}

More generally, we can obtain other $\ell$EAS with the help of a bialgebra projection or with certain linear forms.

\begin{Proposition}
Let $(A,m,\Delta)$ be a bialgebra, not necessarily unitary nor counitary, and $\pi\colon A\longrightarrow A$ be a bialgebra morphism such that $\pi^2=\pi$. For any $a,b\in A$, we define $\Phi\colon A\otimes A\longrightarrow A\otimes A$~by
\begin{align*}
\Phi(a\otimes b)=(m\otimes \pi)\circ (\operatorname{Id}_A\otimes \tau)\circ (\Delta\otimes \operatorname{Id}_A)(a\otimes b)=\sum a^{(1)}b\otimes \pi\bigl(a^{(2)}\bigr).
\end{align*}
 Then $(A,\Phi)$ is an $\ell$EAS.
\end{Proposition}

\begin{proof}
We define $\delta=(\operatorname{Id} \otimes \pi)\circ \Delta$. Then $(A,m,\delta)$ is a bialgebra. Note that it is not counitary, except if $(A,\Delta)$ is counitary and $\pi=\operatorname{Id}_A$. We can then apply Proposition \ref{prop5.1} to $(A,m,\delta)$.
\end{proof}

\begin{Example}
Let $(\Omega,\star)$ be a semigroup and $\pi\colon\Omega\longrightarrow \Omega$ be a semigroup morphism such that~${\pi^2=\pi}$. We take $A=\mathbb{K}\Omega$, with its usual bialgebra structure. Then in $\leas(A,m,\Delta)$, for any $\alpha,\beta \in \Omega$,
$\Phi(\alpha \otimes \beta)=\alpha\star \beta \otimes \pi(\alpha)$.
We recover the linearization of $\eas(\Omega,\star,\pi)$.
\end{Example}

\begin{Proposition}
Let $(A,\Delta)$ be a coalgebra, not necessarily counitary, and $f\in A^*$ such that $(f\otimes f)\circ \Delta=f$. We put, for any $a,b\in A$,
\[\Phi(a\otimes b)=\sum f\bigl(a^{(1)}\bigr)b\otimes a^{(2)}.\]
Then $(A,\Phi)$ is an $\ell$CEDS.
\end{Proposition}

\begin{proof}
We define a product on $A$ by $a\star b=f(a)b$. It is associative. Moreover, for any $a,b\in A$, as $(f\otimes f)\circ \Delta=f$,
\begin{align*}
\Delta(a\star b)&=f(a)\sum b^{(1)}\otimes b^{(2)}=\sum f\bigl(a^{(1)}\bigr)f\bigl(a^{(2)}\bigr)\sum b^{(1)}\otimes b^{(2)}\\
&=\sum\sum f\bigl(a^{(1)}\bigr)b^{(1)}\otimes f\bigl(a^{(2)}\bigr)b^{(2)}=\Delta(a)\star \Delta(b),
\end{align*}
so $(A,\star,\Delta)$ is a bialgebra, and $(A,\Phi)=\leas(A,\star,\Delta)$. Moreover, for any $a,b,c\in A$,
\[a\star b\star c=f(a)f(b)c=f(b)f(c)a=b\star a\star c.\]
By Proposition \ref{prop5.2}, $(A,\Phi)$ is an $\ell$CEDS. \end{proof}

\begin{Example}
Let $\Omega$ be a set, $A=\mathbb{K}\Omega$ be the associated coalgebra (where any $\alpha \in \Omega$ is a~group-like element), and $\Omega'\subseteq \Omega$ be any set. We define the linear form $f\colon A\longrightarrow \mathbb{K}$ by
\begin{align*}
 \forall \alpha \in \Omega, f(\alpha) =\begin{cases}
1& \text{if }\alpha \in \Omega',\\
0 & \text{otherwise}.
\end{cases} \end{align*}
For any $\alpha \in \Omega$, $(f\otimes f)\circ \Delta(\alpha)=f(\alpha)^2=f(\alpha)$, so we obtain an $\ell$CEDS such that for any $\alpha,\beta \in \Omega$,
\[\Phi(\alpha \otimes \beta)=\begin{cases}
\beta\otimes \alpha&\text{if }\alpha \in \Omega',\\
0 & \text{otherwise}.
\end{cases}\] \end{Example}


\subsection[A functor from Hopf algebras to lEAS]{A functor from Hopf algebras to $\boldsymbol{\ell}$EAS}

\begin{Proposition}\label{prop5.7}
Let $(A,m,\Delta)$ be a Hopf algebra, of antipode $S$. For any $a,b\in A$, we define $\Phi\colon A\otimes A\longrightarrow A\otimes A$ by
\begin{gather*}
\Phi(a\otimes b)=(\operatorname{Id}_A\otimes m)\circ (\operatorname{Id}_A\otimes S\otimes \operatorname{Id}_A)\circ (\Delta \otimes \operatorname{Id})\circ \tau(a\otimes b)=\sum b^{(1)}\otimes S\bigl(b^{(2)}\bigr)a.
\end{gather*}
Then $(A,\Phi)$ is an $\ell$EAS, denoted by $\leas'(A,m,\Delta)$. It is non-degenerate, and $\bigl(A,\Phi^{-1}\bigr)=\leas(A,m,\Delta^{\rm op})$.
\end{Proposition}

\begin{proof}
Let $a,b,c\in A$. Then
\begin{gather*}
(\operatorname{Id} \otimes \Phi)\circ(\Phi\otimes \operatorname{Id})\circ (\operatorname{Id} \otimes \Phi)(a\otimes b\otimes c)\\
\qquad=\sum \sum c^{(1)}\otimes S\bigl(c^{(3)}\bigr)^{(1)}b^{(1)}\otimes S\bigl(S\bigl(c^{(3)}\bigr)^{(2)}b^{(2)}\bigr)S\bigl(c^{(2)}\bigr)a\\
\qquad=\sum \sum c^{(1)}\otimes S\bigl(c^{(4)}\bigr)b^{(1)}\otimes S\bigl(S\bigl(c^{(3)}\bigr)b^{(2)}\bigr)S\bigl(c^{(2)}\bigr)a\\
\qquad=\sum \sum c^{(1)}\otimes S\bigl(c^{(4)}\bigr)b^{(1)}\otimes S\bigl(c^{(2)}S\bigl(c^{(3)}\bigr)b^{(2)}\bigr)a\\
\qquad=\sum \sum c^{(1)}\otimes S\bigl(c^{(2)}\bigr)b^{(1)}\otimes S\bigl(b^{(2)}\bigr)a\\
\qquad=(\Phi\otimes \operatorname{Id})\circ (\operatorname{Id} \otimes \tau)\circ (\Phi\otimes \operatorname{Id})(a\otimes b\otimes c),
\end{gather*}
so $(A,\Phi)$ is an $\ell$EAS.

Let $(A,\Psi)=\leas(A,m,\Delta^{\rm op})$: for any $a,b\in A$, $\Psi(a\otimes b)=\sum a^{(2)}b\otimes a^{(1)}$. Then
\begin{gather*}
\Phi\circ \Psi(a \otimes b)=\sum a^{(1)}\otimes S\bigl(a^{(2)}\bigr)a^{(3)}b=a\otimes b,\\
\Psi\circ \Phi(a\otimes b)=\sum b^{(1)}S\bigl(b^{(2)}\bigr)a\otimes b^{(3)}=a\otimes b,
\end{gather*}
so $\Phi$ is bijective, of inverse $\Psi$. \end{proof}

\begin{Example}
Let $(G,\star)$ be a group and let $A=\mathbb{K} G^{\rm op}$ be the Hopf algebra of the opposite of this group. A basis of $\leas'(A,m,\Delta)$ is given by $G$ itself and, for any $\alpha,\beta\in G$,
$\Phi(\alpha \otimes \beta)=\beta \otimes \alpha \star \beta^{-1}$.
We recover in this way the linearization of $\eas'(G,\star)$.
\end{Example}

\begin{Corollary}
Let $(A,m,\Delta)$ be a bialgebra, such that $(A,m,\Delta^{\rm op})$ is a bialgebra. Then $(A,\Phi)=\leas(A,m,\Delta)$ is non-degenerate and $\bigl(A,\Phi^{-1}\bigr)=\leas'(A,m,\Delta^{\rm op})$.
\end{Corollary}

\begin{Proposition}
Let $(A,m,\Delta)$ be a Hopf algebra.
\begin{enumerate}\itemsep=0pt
\item[$(1)$] Then $\leas'(A,m,\Delta)$ is an $\ell$CEDS if and only if $\Delta \circ S=\Delta^{\rm op}\circ S$.
\item[$(2)$] Then $\leas'(A,m,\Delta)$ is a dual $\ell$CEDS if and only if $S\circ m=S\circ m^{\rm op}$.
\end{enumerate}
\end{Proposition}

\begin{proof}
(1) Let $a,b,c\in A$.
\begin{align}
\label{eq34} (\operatorname{Id} \otimes \Phi)\circ (\operatorname{Id} \otimes \tau)\circ (\tau \otimes \operatorname{Id})\circ (\Phi\otimes \operatorname{Id})(b\otimes c\otimes a)&=\sum S\bigl(b^{(3)}\bigr)a\otimes b^{(1)}\otimes S\bigl(b^{(2)}\bigr)c,\\
\nonumber (\tau \otimes \operatorname{Id})\circ (\Phi\otimes \operatorname{Id})\circ (\operatorname{Id} \otimes \Phi)\circ (\operatorname{Id} \otimes \tau)(b\otimes c\otimes a)&=\sum S\bigl(b^{(2)}\bigr)a\otimes b^{(1)}\otimes S\bigl(b^{(3)}\bigr)c.
\end{align}
If $\Delta \circ S=\Delta^{\rm op}\circ S$, then
\begin{align*}
\sum b^{(1)}\otimes S\bigl(b^{(2)}\bigr)\otimes S\bigl(b^{(3)}\bigr)&=\sum \sum b^{(1)}\otimes S\bigl(b^{(2)}\bigr)^{(2)} \otimes S\bigl(b^{(2)}\bigr)^{(1)}\\
&=\sum \sum b^{(1)}\otimes S\bigl(b^{(2)}\bigr)^{(1)} \otimes S\bigl(b^{(2)}\bigr)^{(2)}\\
&=\sum b^{(1)}\otimes S\bigl(b^{(3)}\bigr)\otimes S\bigl(b^{(2)}\bigr),
\end{align*}
which implies that $(A,\Phi)$ is an $\ell$CEDS. Conversely, taking $a=c=1_A$, we obtain, in \eqref{eq34},
\[\sum S\bigl(b^{(3)}\bigr)\otimes b^{(1)}\otimes S\bigl(b^{(2)}\bigr)=\sum S\bigl(b^{(2)}\bigr)\otimes b^{(1)}\otimes S\bigl(b^{(3)}\bigr).\]
Applying $\operatorname{Id} \otimes \varepsilon \otimes \operatorname{Id}$, we obtain
\begin{align*}
\Delta \circ S(b)&=\sum S(b)^{(1)}\otimes S(b)^{(2)}=\sum S\bigl(b^{(2)}\bigr)\otimes S\bigl(b^{(1)}\bigr)\\
&=\sum S\bigl(b^{(1)}\bigr)\otimes S\bigl(b^{(2)}\bigr)=\sum S(b)^{(2)}\otimes S(b)^{(1)}=\Delta^{\rm op}\circ S(b).
\end{align*}

(2) Let $a,b,c\in A$. Then
\begin{gather*}
(\Phi\otimes \operatorname{Id})\circ (\tau \otimes \operatorname{Id})\circ (\operatorname{Id} \otimes \tau)\circ (\operatorname{Id} \otimes \Phi)(b\otimes a\otimes c)=\sum \sum a^{(1)}\otimes S\bigl(a^{(2)}\bigr)S\bigl(c^{(2)}\bigr)b\otimes c^{(1)},\\
\nonumber(\operatorname{Id} \otimes \tau)\circ (\operatorname{Id} \otimes \Phi)\circ (\Phi\otimes \operatorname{Id})\circ (\tau \otimes \operatorname{Id})(b\otimes a\otimes c)=\sum \sum a^{(1)}\otimes S\bigl(c^{(2)}\bigr)S\bigl(a^{(2)}\bigr)b\otimes c^{(1)}.
\end{gather*}
If $S\circ m=S\circ m^{\rm op}$, then $m^{\rm op}\circ (S\otimes S)=m\circ (S\otimes S)$, which implies that $(A,\Phi)$ is a dual $\ell$CEDS. Conversely, taking $b=1$ and applying $\varepsilon \otimes \operatorname{Id} \otimes \varepsilon$, we obtain
$S(a)S(c)=S(c)S(a)$,
so~${m^{\rm op}\circ S=m\circ (S\otimes S)=m^{\rm op}\circ (S\otimes S)=S\circ m}$. \end{proof}

\begin{Remark}
In particular, if $S$ is invertible, then $\leas'(A,m,\Delta)$ is an $\ell$CEDS if and only if $(A,m)$ is commutative; it is a dual $\ell$CEDS if and only if $(A,\Delta)$ is cocommutative.
\end{Remark}

\begin{Proposition}
Let $(A,m,\Delta)$ be a finite-dimensional Hopf algebra. Then \[
\leas'(A,m,\Delta)^* =\leas'(A^*,\Delta^{*,{\rm op}},m^{*,{\rm op}}).\]
\end{Proposition}

\begin{proof}
Let $f,g\in A^*$. For any $a,b\in A$,
\begin{align*}
\Phi^*(f\otimes g)(a\otimes b)&=(f\otimes g)(\Phi(a\otimes b))=\sum (f\otimes g)\bigl(b^{(1)}\otimes S\bigl(b^{(2)}a\bigr)\bigr)\\
&=\sum \sum \bigl(f\otimes g^{(1)}\otimes g^{(2)}\bigr)\bigl(b^{(1)}\otimes S\bigl(b^{(2)}\bigr)\otimes a\bigr)\\
&=\sum \bigl(f\otimes S^*\bigl(g^{(1)}\bigr)\otimes g^{(2)}\bigr)\bigl(b^{(1)}\otimes b^{(2)}\otimes a\bigr)\\
&=\sum \bigl(g^{(2)}\otimes fS^*\bigl(g^{(1)}\bigr)\bigr)(a\otimes b),
\end{align*}
so $\Phi^*(f\otimes g)=\sum g^{(2)}\otimes fS^*\bigl(g^{(1)}\bigr)$, which is the $\ell$EAS attached to the Hopf algebra $(A^*,\Delta^{*,{\rm op}},\allowbreak m^{*,{\rm op}})$, whose antipode is $S^*$. \end{proof}

Recall from \cite{Sweedler1969-2} that a right integral of a Hopf algebra $(A,m,\Delta)$ is a linear map $f\in A^*$ such that for any $\mu\in A^*$,
$(\lambda \otimes \mu)\circ \Delta=\mu(1_A) \lambda$.

\begin{Proposition} \label{prop5.11}
Let $(A,m,\Delta)$ be a Hopf algebra.
\begin{enumerate}\itemsep=0pt
\item[$(1)$] Let $a\in A$. It is a left unit of $\leas'(A,m,\Delta)$ if and only if for any $b\in A$, $S(b)a=\varepsilon(b)a$.
\item[$(2)$] Let $\lambda \in A^*$. It is a left counit of $\leas'(A,m,\Delta)$ if and only if for any $a\in A$, \linebreak $\sum \lambda\bigl(b^{(1)}\bigr) S\bigl(b^{(2)}\bigr)=\lambda(b)1_A$.
In particular, right integrals on $(A,m,\Delta)$ are left counit of $\leas'(A,m,\Delta)$; if $S$ is invertible, then the converse is true.
\end{enumerate}\end{Proposition}

\begin{proof}
(1) Let $a\in A$. Then its a left unit if and only if for any $b\in A$, $\sum b^{(1)}\otimes S\bigl(b^{(2)}\bigr)a=b\otimes a$. Applying $\varepsilon\otimes \operatorname{Id}$, if $a$ is a left unit, for any $b\in B$, $S(b)a=\varepsilon(b)\otimes a$.
Conversely, if this holds, then for any $b\in B$,
\[\Phi(a\otimes b)=\sum b^{(1)}\otimes S\bigl(b^{(1)}\bigr)a=\sum b^{(1)}\otimes \varepsilon\bigl(b^{(2)}\bigr) a=b\otimes a.\]

(2) Let $\lambda \in A^*$. It is a left counit if and only if for any $a,b\in A$,
\[\sum \lambda\bigl(b^{(1)}\bigr)S\bigl(b^{(2)}\bigr)a=a\lambda(b).\]
If $\lambda$ is a left counit, taking $a=1_A$, we obtain that for any $b\in A$,
$\lambda\bigl(b^{(1)}\bigr)S\bigl(b^{(2)}\bigr)=\lambda(b)1_A$. Conversely, if this holds, then for any $a,b\in A$,{\samepage
\[(\lambda\otimes \operatorname{Id})\circ \Phi(a\otimes b)=\sum \lambda\bigl(b^{(1)}\bigr)S\bigl(b^{(2)}\bigr)a=\lambda(b)a=(\operatorname{Id} \otimes \lambda)(a\otimes b),\]
so $\lambda$ is a left counit.}

Let us assume that $\lambda$ is a right integral of $(A,m,\Delta)$. For any $b\in A$, for any $\mu\in A^*$,
\begin{align*}
\sum \lambda\bigl(b^{(1)}\bigr)\mu\bigl(S\bigl(b^{(2)}\bigr)\bigr)&=(\lambda\otimes \mu\circ S)\circ \Delta(b)=\mu\circ S(1_A)\lambda(b)=\mu(1_A)\lambda(b).
\end{align*}
As this holds for any $\mu\in A^*$, $\sum \lambda\bigl(b^{(1)}\bigr)S\bigl(b^{(2)}\bigr)=\lambda(b)1_A$, so $\lambda$ is a right integral. Let us now assume that $S$ is invertible and that $\lambda$ is a left counit.
Let $\nu \in A^*$. For any $b\in A$, if $\mu=\nu \circ S^{-1}$,
\begin{align*}
\sum \lambda\bigl(b^{(1)}\bigr)\nu\bigl(b^{(2)}\bigr)&=\sum \lambda\bigl(b^{(1)}\bigr)\mu\circ S\bigl(b^{(2)}\bigr)=\lambda(b)\mu(1_A)\\
&=\lambda(b)\nu\circ S^{-1}(1_A)=\lambda(b)\nu(1_A).
\end{align*}
So $\lambda$ is a right integral.
\end{proof}


\subsection{From left units and counits to bialgebras}

\begin{Theorem}\label{theo5.12}
Let $(A,\Phi)$ be an $\ell$EAS.
\begin{enumerate}\itemsep=0pt
\item[$(1)$] If $a$ is a special vector of eigenvalue $1$ of $(A,\Phi)$, then $\Delta_a\colon A\longrightarrow A\otimes A$ defined by $\Delta_a(b)=\Phi(b\otimes a)$ is a coassociative coproduct.
\item[$(2)$] If $\varepsilon$ is a special vector of eigenvalue $1$ of $(A,\Phi)^*$, that is to say if $(\varepsilon\otimes \varepsilon)\circ \Phi=\varepsilon\otimes \varepsilon$,
 then $m_\varepsilon\colon A\otimes A\longrightarrow A$ defined by $m_\varepsilon=(\operatorname{Id} \otimes \varepsilon)\circ \Phi$ is an associative product.
\item[$(3)$] If $a$ is a left unit of $(A,\Phi)$ and $\varepsilon$ is a left counit of $(A,\Phi)$ such that $\varepsilon(a)=1$, then $(A,m_\varepsilon,\Delta_a)$ is a bialgebra, with $a$ as a left unit and $\varepsilon$ as a left counit.
Moreover, $(A,\Phi)=\leas(A,m_\varepsilon,\Delta_a)$.
\end{enumerate}\end{Theorem}

\begin{proof}
(1) For any $b\in A$,
\begin{gather*}
(\operatorname{Id} \otimes \Phi)\circ (\Phi\otimes \operatorname{Id})\circ (\operatorname{Id} \otimes \Phi)(b\otimes a\otimes a)=(\operatorname{Id} \otimes \Phi)\circ (\Phi\otimes \operatorname{Id})(b\otimes a\otimes a)\\
\phantom{(\operatorname{Id} \otimes \Phi)\circ (\Phi\otimes \operatorname{Id})\circ (\operatorname{Id} \otimes \Phi)(b\otimes a\otimes a)}{}=(\operatorname{Id} \otimes \Phi)(\Delta_a(b)\otimes a)=(\operatorname{Id} \otimes \Delta_a)\circ \Delta_a(b),\\
(\Phi\otimes \operatorname{Id})\circ (\operatorname{Id} \otimes \tau)\circ (\Phi\otimes \operatorname{Id})(b\otimes a\otimes a)=(\Phi\otimes \operatorname{Id})\circ (\operatorname{Id} \otimes \tau)(\Delta_a(b)\otimes a)\\
\phantom{(\Phi\otimes \operatorname{Id})\circ (\operatorname{Id} \otimes \tau)\circ (\Phi\otimes \operatorname{Id})(b\otimes a\otimes a)}{}=(\Delta_a\otimes \operatorname{Id})\circ \Delta_a(b).
\end{gather*}
Hence, $\Delta_a$ is coassociative.

(2) We obtain, as $\varepsilon$ is a special vector of eigenvalue $1$ of $(A,\Phi)^*$,
\begin{gather*}
(\operatorname{Id} \otimes \varepsilon\otimes \varepsilon)\circ (\operatorname{Id} \otimes \Phi)\circ (\Phi\otimes \operatorname{Id})\circ (\operatorname{Id} \otimes \Phi)=(\operatorname{Id} \otimes \varepsilon\otimes \varepsilon)\circ (\Phi\otimes \operatorname{Id})\circ (\operatorname{Id} \otimes \Phi)\\
\phantom{(\operatorname{Id} \otimes \varepsilon\otimes \varepsilon)\circ (\operatorname{Id} \otimes \Phi)\circ (\Phi\otimes \operatorname{Id})\circ (\operatorname{Id} \otimes \Phi)}{}=(\operatorname{Id} \otimes \varepsilon)\circ \Phi\circ ((\operatorname{Id} \otimes \varepsilon)\circ \Phi))\\
\phantom{(\operatorname{Id} \otimes \varepsilon\otimes \varepsilon)\circ (\operatorname{Id} \otimes \Phi)\circ (\Phi\otimes \operatorname{Id})\circ (\operatorname{Id} \otimes \Phi)}{}=m_\varepsilon\circ(\operatorname{Id} \otimes m_\varepsilon),\\
(\operatorname{Id} \otimes \varepsilon\otimes \varepsilon)\circ (\Phi\otimes \operatorname{Id})\circ (\operatorname{Id} \otimes \tau)\circ (\Phi\otimes \operatorname{Id})=(\operatorname{Id} \otimes \varepsilon )\circ \Phi\circ ((\operatorname{Id} \otimes \varepsilon )\otimes \operatorname{Id})\\
\phantom{(\operatorname{Id} \otimes \varepsilon\otimes \varepsilon)\circ (\Phi\otimes \operatorname{Id})\circ (\operatorname{Id} \otimes \tau)\circ (\Phi\otimes \operatorname{Id})}{}=m_\varepsilon\circ(m_\varepsilon\otimes \operatorname{Id}).
\end{gather*}
As a consequence, $m_\varepsilon$ is associative.

(3) As $a$ is a left unit, it is a special vector of eigenvalue $1$ of $(A,\Phi)$, so $\Delta_a$ is coassociative. Moreover, for any $b\in A$,
\[(\varepsilon\otimes \operatorname{Id})\circ \Delta_a(b)=(\varepsilon\otimes \operatorname{Id})\circ \Phi(b\otimes a)=(\operatorname{Id} \otimes \varepsilon)(b\otimes a)=b\varepsilon(a)=b,\]
so $\varepsilon$ is a left counit of $\Delta_a$. As $\varepsilon$ is a left counit, it is a special vector of eigenvalue $1$ of $(A,\Phi)^*$, so $m_\varepsilon$ is associative. Moreover, for any $b\in A$,{\samepage
\[m_\varepsilon(a\otimes b)=(\operatorname{Id} \otimes \varepsilon)\circ \Phi(a\otimes b)=(\operatorname{Id} \otimes \varepsilon)(b\otimes a)=b\varepsilon(a)=b,\]
so $a$ is a left unit of $m_\varepsilon$.}

Let $b_1,b_2\in A$.
\begin{align*}
\Delta_a(b_1b_2)={}&(\operatorname{Id} \otimes \varepsilon\otimes \operatorname{Id} \otimes \varepsilon)\circ (\Phi\otimes \Phi)\circ (\operatorname{Id} \otimes \tau\otimes \operatorname{Id})\circ (\Phi\otimes \Phi)\\
&\circ (\operatorname{Id} \otimes \tau\otimes \operatorname{Id})(b_1\otimes b_2\otimes a\otimes a)\\
={}&(\operatorname{Id} \otimes \varepsilon\otimes \operatorname{Id} \otimes \varepsilon)\circ (\operatorname{Id} \otimes \operatorname{Id} \otimes \Phi)\circ (\Phi\otimes \operatorname{Id} \otimes \operatorname{Id})\circ (\operatorname{Id} \otimes \tau\otimes \operatorname{Id})\circ (\Phi\otimes \operatorname{Id} \otimes \operatorname{Id})\\
&\circ(\operatorname{Id} \otimes \operatorname{Id} \otimes \Phi)\circ (\operatorname{Id} \otimes \tau\otimes \operatorname{Id})(b_1\otimes b_2\otimes a\otimes a)\\
={}&(\operatorname{Id} \otimes \varepsilon\otimes \operatorname{Id} \otimes \varepsilon)\circ (\operatorname{Id} \otimes \operatorname{Id} \otimes \Phi)\circ (\operatorname{Id}\otimes \Phi \otimes \operatorname{Id})\circ (\Phi \otimes \operatorname{Id}\otimes \operatorname{Id})\circ (\operatorname{Id}\otimes \Phi\otimes \operatorname{Id})\\
&\circ(\operatorname{Id} \otimes \operatorname{Id} \otimes \Phi)\circ (\operatorname{Id} \otimes \tau\otimes \operatorname{Id})(b_1\otimes b_2\otimes a\otimes a)\\
={}&(\operatorname{Id} \otimes \operatorname{Id} \otimes \varepsilon)\circ (\operatorname{Id} \otimes \Phi)\circ (\operatorname{Id}\otimes ((\varepsilon \otimes \operatorname{Id})\circ\Phi) \otimes \operatorname{Id})\circ (\Phi \otimes \operatorname{Id}\otimes \operatorname{Id})\\
&\circ (\operatorname{Id}\otimes \Phi\otimes \operatorname{Id})\circ(\operatorname{Id} \otimes \operatorname{Id} \otimes \Phi)\circ (\operatorname{Id} \otimes \tau\otimes \operatorname{Id})(b_1\otimes b_2\otimes a\otimes a)\\
={}&(\operatorname{Id} \otimes \operatorname{Id} \otimes \varepsilon)\circ (\operatorname{Id} \otimes \Phi)\circ (\operatorname{Id}\otimes \operatorname{Id} \otimes \varepsilon \otimes \operatorname{Id})\circ (\Phi \otimes \operatorname{Id}\otimes \operatorname{Id})\\
&\circ (\operatorname{Id}\otimes \Phi\otimes \operatorname{Id})\circ(\operatorname{Id} \otimes \operatorname{Id} \otimes \Phi)\circ (\operatorname{Id} \otimes \tau\otimes \operatorname{Id})(b_1\otimes b_2\otimes a\otimes a)\\
={}&(\operatorname{Id} \otimes \operatorname{Id} \otimes \varepsilon)\circ (\operatorname{Id} \otimes \Phi)\circ (\Phi \otimes \operatorname{Id}\otimes \operatorname{Id})\circ (\operatorname{Id}\otimes (\operatorname{Id}\otimes \varepsilon)\circ\Phi\otimes \operatorname{Id})\\
&\circ(\operatorname{Id} \otimes \operatorname{Id} \otimes \Phi)(b_1\otimes a\otimes b_2\otimes a)\\
={}&(\operatorname{Id} \otimes \operatorname{Id} \otimes \varepsilon)\circ (\operatorname{Id} \otimes \Phi)\circ (\Phi \otimes \operatorname{Id})\circ(\operatorname{Id} \otimes \Phi)(b_1\otimes b_2\otimes a)
\\
={}&\Phi\circ((\operatorname{Id} \otimes \varepsilon)\circ \Phi \otimes \operatorname{Id})(b_1\otimes b_2\otimes a)\\
={}&\Phi(m_\varepsilon(b_1\otimes b_2)\otimes a)
=\Delta_a(m_\varepsilon(b_1\otimes b_2)).
\end{align*}
So $(A,m_\varepsilon,\Delta_a)$ is a bialgebra. Let $(A,\Psi)=\leas(A,m_\varepsilon,\Delta_a)$. For any $b_1,b_2\in A$,
\begin{align*}
\Psi(b_1\otimes b_2)&=(\operatorname{Id} \otimes \varepsilon\otimes \operatorname{Id})\circ (\Phi\otimes \operatorname{Id})\circ (\operatorname{Id} \otimes \tau)\circ (\Phi\otimes \operatorname{Id})(b_1\otimes a\otimes b_2)\\
&=(\operatorname{Id} \otimes \varepsilon\otimes \operatorname{Id})\circ (\operatorname{Id}\otimes \Phi)\circ (\Phi \otimes \operatorname{Id})\circ (\operatorname{Id}\otimes \Phi)(b_1\otimes a\otimes b_2)\\
&=(\operatorname{Id} \otimes \operatorname{Id} \otimes \varepsilon)\circ (\Phi\otimes \operatorname{Id})(b_1\otimes b_2\otimes a)=\Phi(b_1\otimes b_2)\varepsilon(a)=\Phi(b_1\otimes b_2).
\end{align*}
Therefore, $(A,\Phi)=\leas(A,m_\varepsilon,\Delta_a)$. \end{proof}

\begin{Example}
This can be applied for $\ell$EAS $M_{16}$, $M_{17}$ and $M_{18}$ of Example \ref{ex4.1}.
\begin{itemize}
\item For $M_{16}$, taking $a=x$ and $\varepsilon=x^*+y^*$, we obtain
\begin{alignat*}{3}
&\Delta_a(x)=x\otimes x,\qquad&&\Delta_a(y)=y\otimes y,&\\
&m_\varepsilon(x\otimes x)=x,\qquad&& m_\varepsilon(x\otimes y)=y,&\\
&m_\varepsilon(y\otimes x)=y,\qquad&& m_\varepsilon(y\otimes y)=y.&
\end{alignat*}
This is the bialgebra of the semigroup $(\mathbb{Z}/2\mathbb{Z},\times)$, with $x=\overline{1}$ and $y=\overline{0}$:
we recover the linearization of $\mathbf{C3}$.
\item For $M_{17}$, taking $a=x$ and $\varepsilon=x^*+y^*$, we obtain
\begin{alignat*}{3}
&\Delta_a(x)=x\otimes x,\qquad&&\Delta_a(y)=x\otimes x-x\otimes y-y\otimes x+2y\otimes y,&\\
&m_\varepsilon(x\otimes x)=x,\qquad&& m_\varepsilon(x\otimes y)=y,&\\
&m_\varepsilon(y\otimes x)=x,\qquad&& m_\varepsilon(y\otimes y)=y.&
\end{alignat*}
Putting $y'=-x+2y$,\footnote{If the characteristic of the base field $\mathbb{K}$ is not 2.} we obtain
\begin{alignat*}{3}
&\Delta_a(x)=x\otimes x,\qquad&&\Delta_a(y')=y'\otimes y',&\\
&m_\varepsilon(x\otimes x)=x,\qquad&& m_\varepsilon(x\otimes y')=y',&\\
&m_\varepsilon(y'\otimes x)=y',\qquad&& m_\varepsilon(y'\otimes y')=x.&
\end{alignat*}
This is the bialgebra of the semigroup $(\mathbb{Z}/2\mathbb{Z},+)$, with $x=\overline{0}$ and $y=\overline{1}$:
we recover the linearization of $\mathbf{H2}$.
\item For $M_{18}$, we can take any $a\in A$ and any $\varepsilon\in A^*$ such that $\varepsilon(a)=1$. For any $b,c\in A$,
$
\Delta_a(b)=a\otimes b$, $ m_\varepsilon(b\otimes c)=\varepsilon(b)c$.
\end{itemize}\end{Example}


\subsection{Applications to non-degenerate finite CEDS}

From Proposition \ref{prop4.5}, we have the following.

\begin{Proposition}
Let $(\Omega,\rightarrow,\triangleright)$ be a non-degenerate finite CEDS, which we write following Theorem {\rm\ref{theo3.14}} under the form
$(\eas(\Omega_1,*) \rtimes_\succ \eas'(\Omega_2,\star))\times \eas(\Omega_3)$.
Let $g,h\colon\Omega_3\longrightarrow \mathbb{K}$ be two maps such that
\[\sum_{\alpha_3\in \Omega_3} g(\alpha_3)h(\alpha_3)=1.\]
We define a product and a coproduct on $\mathbb{K}\Omega$, putting, for any $(\alpha_1,\alpha_2,\alpha_3)$, $(\beta_1,\beta_2,\beta_3)\in \Omega$,
\begin{gather*}
(\alpha_1,\alpha_2,\alpha_3)\cdot(\beta_1,\beta_2,\beta_3)=\delta_{\alpha_2,\beta_2}g(\alpha_3)(\alpha_1*\beta_1,\beta_2,\beta_3),\\
\Delta(\alpha_1,\alpha_2,\alpha_3)=\sum_{(\beta_2,\beta_3)\in \Omega_2\times \Omega_3}h(\beta_3)(\alpha_1,\beta_2,\beta_3)\otimes \bigl(\beta_2\succ \alpha_1,\alpha_2\star \beta_2^{-1},\alpha_3\bigr).
\end{gather*}
Then $(\mathbb{K}\Omega,\cdot,\Delta)$ is a bialgebra and the linearization of $\Omega$ is $\leas(\mathbb{K}\Omega,\cdot,\Delta)$.
\end{Proposition}

\begin{proof}
By Proposition \ref{prop4.5}, the following is a left unit of $\mathbb{K}\Omega$:
\[a=\sum_{(\alpha_2,\alpha_3)\in \Omega_2\times \Omega_3} h(\alpha_3)(e_1,\alpha_2,\alpha_3),\]
and the following map is a left counit of $\mathbb{K}\Omega$:
\[\varepsilon\colon \ \begin{cases}
\mathbb{K}\Omega \longrightarrow \mathbb{K},\\
(\alpha_1,\alpha_2,\alpha_3) \longrightarrow \delta_{\alpha_2,e_2}g(\alpha_3).
\end{cases}\]
By hypothesis, $\varepsilon(a)=1$. The result comes from a direct application of Theorem~\ref{theo5.12}.
\end{proof}

Similarly, we have the following.

\begin{Proposition}
Let $(\Omega,\rightarrow,\triangleright)$ be a non-degenerate finite dual CEDS, which we write following Corollary {\rm\ref{cor3.15}} under the form
$(\eas(\Omega_2,\star) \ltimes_\prec (\eas'(\Omega_1,*))\times \eas(\Omega_3)$,
Let $g,h\colon\Omega_3\longrightarrow \mathbb{K}$ be two maps such that
\[\sum_{\alpha_3\in \Omega_3} g(\alpha_3)h(\alpha_3)=1.\]
We define a product and a coproduct on $\mathbb{K}\Omega$, putting, for any $(\alpha_1,\alpha_2,\alpha_3)$, $(\beta_1,\beta_2,\beta_3)\in \Omega$,
\begin{gather*}
(\alpha_1,\alpha_2,\alpha_3)\cdot(\beta_1,\beta_2,\beta_3)=\delta_{\alpha_1,\beta_1}g(\alpha_3)(\alpha_2*\beta_2,\beta_1\prec \alpha_2,\beta_3),\\
\Delta(\alpha_2,\alpha_1,\alpha_3)=\sum_{(\beta_1,\beta_3)\in \Omega_1\times \Omega_3}h(\beta_3)(\alpha_2,\beta_1\prec \alpha_2,\beta_3)\otimes \bigl(\alpha_2,\alpha_1\star\bigl(\beta_1^{-1}\prec \alpha_2^{-1}\bigr),\alpha_3\bigr).
\end{gather*}
Then $(\mathbb{K}\Omega,\cdot,\Delta)$ is a bialgebra and the linearization of $\Omega$ is $\leas(\mathbb{K}\Omega,\cdot,\Delta)$.
\end{Proposition}

\subsection{Applications to Hopf algebras of groups}

In all this paragraph, $G$ is a group. We denote by $\mathbb{K} G$ the associated Hopf algebra. If $G$ is finite, we denote by $\mathbb{K}^G$ the Hopf algebra of functions over $G$, with its basis $(\delta_g)_{g\in G}$, dual of the basis~$G$ of $\mathbb{K} G$.

\begin{Corollary}\label{cor5.15}
If $G$ is finite, then $\leas'(\mathbb{K} G)$ is isomorphic to $\leas\bigl(\mathbb{K}^G\bigr)$, and $\leas'\bigl(\mathbb{K}^G\bigr)$ is isomorphic to $\leas(\mathbb{K} G^{\rm op})$.
\end{Corollary}

\begin{proof}
As $G$ is finite, $a= \sum_{g\in G} g$ is a right integral of $\mathbb{K}^G$, so is a left unit of $\leas'(\mathbb{K} G)$.
If $e_G$ is the unit of the group $G$, then $\varepsilon=\delta_{e_G}$ is a right integral of $\mathbb{K} G$, so is a left counit of~$\leas'(\mathbb{K} G)$. As $\varepsilon(a)=1$, $\leas'(\mathbb{K} G)=\leas(\mathbb{K} G,m_\varepsilon,\Delta_a)$. For any $g,h\in G$,
\[m_\varepsilon(g\otimes h)=(\operatorname{Id} \otimes \delta_{e_G})\circ \Phi(g\otimes h)=h \delta_{e_G}\bigl(h^{-1}g\bigr)=\delta_{g,h} h.\]
For any $g\in G$,
\[\Delta_a(g)=\sum_{h\in G}\Phi(g\otimes h)=\sum_{h\in G}h\otimes h^{-1}g=\sum_{\substack{g_1,g_2\in G,\\ g_1g_2=g}}g_1\otimes g_2.\]
So $(\mathbb{K} G,m_\varepsilon,\Delta_a)$ is isomorphic to $\mathbb{K}^G$, via the map sending $g$ to $\delta_g$, for any $g\in G$.

By duality, $a$ is a left counit of $\leas'\bigl(\mathbb{K}^G\bigr)$ and $\varepsilon$ is a left unit of $\leas'\bigl(\mathbb{K}^G\bigr)$. For any $g,h\in G$,
\[m_a(\delta_g \otimes \delta_h)=(\operatorname{Id} \otimes a)\circ \Phi(g\otimes h)=\sum_{\substack{h_1,h_2\in G,\\ h_1h_2=h}} \delta_{h_1}\otimes \delta_{h_2^{-1}}\delta_g(a)=\delta_{hg}. \]
For any $g\in G$,
\[\Delta_\varepsilon(\delta_g)=\Phi(\delta_g \otimes \delta_{e_G})=\sum_{h\in G}\delta_h \otimes \delta_h \delta_g=\sum_{h\in G}\delta_h \otimes \delta_{g,h} \delta_h=\delta_g \otimes \delta_g.\]
So $\bigl(\mathbb{K}^G,m_a,\Delta_\varepsilon\bigr)$ is isomorphic to $\mathbb{K} G^{\rm op}$ via the map sending $\delta_g$ to $g$, for any $g\in G$. \end{proof}

\begin{Proposition} \label{prop5.16}\quad
\begin{enumerate}\itemsep=0pt
\item[$(1)$] The nonzero special vectors of eigenvalue $1$ of $\leas(\mathbb{K} G)$ and of $\leas'(\mathbb{K} G)$ are the elements
$\lambda \sum_{\alpha \in H} \alpha$,
where $\lambda$ is a nonzero scalar and $H$ is a subgroup of $G$.
\item[$(2)$] If $G$ is finite, the nonzero special vectors of eigenvalue $1$ of $\leas\bigl(\mathbb{K}^G\bigr)$ and of $\leas'\bigl(\mathbb{K}^G\bigr)$ are the elements
$\lambda \sum_{\alpha \in H} \delta_\alpha$,
where $\lambda$ is a nonzero scalar and $H$ is a subgroup of $G$.
\end{enumerate} \end{Proposition}

\begin{proof}
Any $a\in A$ can be written under the form $ a=\sum_{\alpha \in G} \lambda_\alpha \alpha$. Then
\begin{align*}
&a\text{ is a special vector of eigenvalue 1 of }\leas(\mathbb{K} G)\\
&\Longleftrightarrow \sum_{\alpha,\beta \in G} a_\alpha a_\beta \alpha \otimes \beta=\sum_{\alpha,\beta\in G}a_\alpha a_\beta \alpha \beta \otimes \alpha\\
&\Longleftrightarrow\sum_{\alpha,\beta \in G} a_\alpha a_\beta \alpha \otimes \beta=\sum_{\alpha,\beta\in G}a_\beta a_{\beta^{-1}\alpha} \alpha \beta \otimes \alpha\\
&\Longleftrightarrow\forall \alpha,\beta \in G, a_\beta(a_\alpha-a_{\beta^{-1}\alpha})=0.
\end{align*}
Let $a$ be a nonzero special vector of eigenvalue 1 of $\leas(\mathbb{K} G)$. Let us put $a_{1_G}=\lambda$ and $H=\{\alpha \in G, a_\alpha \neq 0\}$. Let $\alpha=\beta \in H$.
As $a_\beta\neq 0$, we obtain $a_\alpha=a_{1_G}=\lambda$, so $1_G\in H$ and~${\lambda \neq 0}$. For any $\beta \in H$, taking $\alpha=1_G$, we~obtain $a_{\beta^{-1}}=\lambda$, so $\beta^{-1}\in H$.
If $\alpha,\beta \in H$, we~obtain that $a_{\beta^{-1}\alpha}=a_\alpha\neq 0$, so $\beta^{-1}\alpha \in H$. Hence, $H$ is a subgroup and $ a=\lambda \sum_{\alpha \in H}\alpha$,
\begin{align*}
&a\text{ is a special vector of eigenvalue 1 of }\leas'(\mathbb{K} G)\\
&\Longleftrightarrow \sum_{\alpha,\beta \in G} a_\alpha a_\beta \alpha \otimes \beta=\sum_{\alpha,\beta\in G}a_\alpha a_\beta \beta \otimes \beta^{-1}\alpha\\
&\Longleftrightarrow\sum_{\alpha,\beta \in G} a_\alpha a_\beta \alpha \otimes \beta=\sum_{\alpha,\beta\in G}a_\beta a_{\alpha\beta} \alpha \beta \otimes \alpha\\
&\Longleftrightarrow\forall \alpha,\beta \in G, a_\alpha(a_\beta-a_{\alpha\beta})=0.
\end{align*}
Let $a$ be a nonzero special vector of eigenvalue 1 of $\leas'(\mathbb{K} G)$. Let us put $a_{1_G}=\lambda$ and $H=\{\alpha \in G, a_\alpha \neq 0\}$. Let $\alpha=\beta \in H$.
If $\alpha \in H$, for $\beta=1_G$, we obtain $a_{1_G}=a_\alpha=\lambda$, so $1_G\in H$ and $\lambda \neq 0$; for $\beta=\alpha^{-1}$, we obtain $a_{\alpha^{-1}}=a_{1_G}=\lambda\neq 0$, so $\alpha^{-1}\in G$.
If $\alpha,\beta \in H$, we obtain that $a_{\alpha\beta}=a_\beta\neq 0$, so $\alpha \beta\in H$. Hence, $H$ is a subgroup and $ a=\lambda \sum_{\alpha \in H}\alpha$.

Let $f\in \mathbb{K}^G$. We put $f(\alpha)=a_\alpha$ for any $\alpha \in G$,
\begin{align*}
&f\text{ is a special vector of eigenvalue 1 of }\leas\bigl(\mathbb{K}^G\bigr)\\
&\Longleftrightarrow\forall \alpha,\beta \in G, a_\alpha a_\beta=a_{\alpha\beta}\\
&\Longleftrightarrow\forall \alpha,\beta \in G, a_\alpha(a_\beta-a_{\alpha\beta})=0;\\
&f\text{ is a special vector of eigenvalue 1 of }\leas'\bigl(\mathbb{K}^G\bigr)\\
&\Longleftrightarrow\forall \alpha,\beta \in G, a_\alpha a_\beta=a_\beta a_{\beta^{-1}\alpha}\\
&\Longleftrightarrow\forall \alpha,\beta \in G, a_\beta(a_\alpha-a_{\beta^{-1}\alpha})=0.
\end{align*}
The conclusion is the same as for $\mathbb{K} G$. \end{proof}

\begin{Remark}\quad
\begin{enumerate}\itemsep=0pt
\item[(1)] From Proposition \ref{prop5.4}, the left units of $\leas(\mathbb{K} G)$ are the multiples of $e_G$, and its left counits are the multiples of its counit.
If $G$ is finite, the left units of $\leas\bigl(\mathbb{K}^G\bigr)$ are the multiple of $ \sum_{g\in G} g$,and its left counits are the multiples of $e_G$.
\item[(2)] From Proposition \ref{prop5.11}, it is not difficult to show that if $G$ is finite, the left units of $\leas'(\mathbb{K} G)$ are the multiples of $ \sum_{g\in G} g$; if $G$ is not finite, $\leas'(\mathbb{K} G)$ has no nonzero left unit.
The left counits of $\leas'(\mathbb{K} G)$ are the multiples of $\delta_{e_G}$.
By duality, if $G$ is finite, the left units of $\leas'\bigl(\mathbb{K}^G\bigr)$ are the multiples of $\delta_{e_G}$ and its left counits are the multiples of $ \sum_{g\in G} g$.
\end{enumerate}\end{Remark}

\subsection*{Acknowledgements}

The author thanks the anonymous referees for their useful and constructive comments on the first version of this text. The author acknowledges support from the grant ANR-20-CE40-0007 \emph{Combinatoire Alg\'e\-brique, R\'esurgence, Probabilit\'es Libres et Op\'erades}.

\pdfbookmark[1]{References}{ref}
\begin{thebibliography}{99}
\footnotesize\itemsep=0pt

\bibitem{Abe1980}
Abe E., Hopf algebras, \textit{Cambridge Tracts in Math.}, Vol.~74, Cambridge
 University Press, Cambridge, 1980.

\bibitem{Bruned2019}
Bruned Y., Hairer M., Zambotti L., Algebraic renormalisation of regularity
 structures, \href{https://doi.org/10.1007/s00222-018-0841-x}{\textit{Invent.
 Math.}} \textbf{215} (2019), 1039--1156,
 \href{http://arxiv.org/abs/1610.08468}{arXiv:1610.08468}.

\bibitem{Bruned2023}
Bruned Y., Katsetsiadis F., Post-{L}ie algebras in regularity structures,
 \href{https://doi.org/10.1017/fms.2023.93}{\textit{Forum Math. Sigma}}
 \textbf{11} (2023), e98, 20~pages,
 \href{http://arxiv.org/abs/2208.00514}{arXiv:2208.00514}.

\bibitem{Bruned2023-1}
Bruned Y., Manchon D., Algebraic deformation for~({S}){PDE}s,
 \href{https://doi.org/10.2969/jmsj/88028802}{\textit{J.~Math. Soc. Japan}}
 \textbf{75} (2023), 485--526,
 \href{http://arxiv.org/abs/2011.05907}{arXiv:2011.05907}.

\bibitem{Cartier2021}
Cartier P., Patras F., Classical {H}opf algebras and their applications,
 \textit{Algebr. Appl.}, Vol.~29,
 \href{https://doi.org/10.1007/978-3-030-77845-3}{Springer}, Cham, 2021.

\bibitem{Ebrahimi-Fard2007}
Ebrahimi-Fard K., Gracia-Bond\'{\i}a J.M., Patras F., A {L}ie theoretic
 approach to renormalization,
 \href{https://doi.org/10.1007/s00220-007-0346-8}{\textit{Comm. Math. Phys.}}
 \textbf{276} (2007), 519--549,
 \href{http://arxiv.org/abs/hep-th/0609035}{arXiv:hep-th/0609035}.

\bibitem{Foissy56}
Foissy L., Generalized associative algebras,
 \href{http://arxiv.org/abs/2104.00908}{arXiv:2104.00908}.

\bibitem{Foissy48}
Foissy L., Generalized prelie and permutative algebras,
 \href{http://arxiv.org/abs/2104.00909}{arXiv:2104.00909}.

\bibitem{Foissy47}
Foissy L., Algebraic structures on typed decorated rooted trees,
 \href{https://doi.org/10.3842/SIGMA.2021.086}{\textit{SIGMA}} \textbf{17}
 (2021), 086, 28~pages,
 \href{http://arxiv.org/abs/1811.07572}{arXiv:1811.07572}.

\bibitem{Foissy49}
Foissy L., Typed binary trees and generalized dendrifom algebras,
 \href{https://doi.org/10.1016/j.jalgebra.2021.06.025}{\textit{J.~Algebra}}
 \textbf{586} (2021), 1--61,
 \href{http://arxiv.org/abs/2002.12120}{arXiv:2002.12120}.

\bibitem{Foissy50}
Foissy L., {M}anchon D., {Z}hang Y., Families of algebraic structures,
 \href{http://arxiv.org/abs/2005.05116}{arXiv:2005.05116}.

\bibitem{Guo2009}
Guo L., Operated semigroups, {M}otzkin paths and rooted trees,
 \href{https://doi.org/10.1007/s10801-007-0119-7}{\textit{J.~Algebraic
 Combin.}} \textbf{29} (2009), 35--62,
 \href{http://arxiv.org/abs/0710.0429}{arXiv:0710.0429}.

\bibitem{Sweedler1969}
Sweedler M.E., Hopf algebras, \textit{Math. Lect. Note Ser.}, W.A.~Benjamin, Inc., New
 York, 1969.

\bibitem{Sweedler1969-2}
Sweedler M.E., Integrals for {H}opf algebras,
 \href{https://doi.org/10.2307/1970672}{\textit{Ann. of Math.}} \textbf{89}
 (1969), 323--335.

\bibitem{Zhang2019}
Zhang Y., Gao X., Free {R}ota--{B}axter family algebras and (tri)dendriform
 family algebras,
 \href{https://doi.org/10.2140/pjm.2019.301.741}{\textit{Pacific~J. Math.}}
 \textbf{301} (2019), 741--766.

\bibitem{Zhang2020}
Zhang Y., Gao X., Guo L., Matching {R}ota--{B}axter algebras, matching
 dendriform algebras and matching pre-{L}ie algebras,
 \href{https://doi.org/10.1016/j.jalgebra.2020.02.011}{\textit{J.~Algebra}}
 \textbf{552} (2020), 134--170,
 \href{http://arxiv.org/abs/1909.10577}{arXiv:1909.10577}.

\bibitem{Zhang2020-3}
Zhang Y., Gao X., Manchon D., Free (tri)dendriform family algebras,
 \href{https://doi.org/10.1016/j.jalgebra.2019.11.027}{\textit{J.~Algebra}}
 \textbf{547} (2020), 456--493,
 \href{http://arxiv.org/abs/1909.08946}{arXiv:1909.08946}.

\bibitem{Manchon2020}
Zhang Y., Manchon D., Free pre-{L}ie family algebras,
 \href{https://doi.org/10.4171/aihpd/162}{\textit{Ann. Inst. Henri
 Poincar\'e~D}} \textbf{11} (2024), 331--361,
 \href{http://arxiv.org/abs/2003.00917}{arXiv:2003.00917}.

\end{thebibliography}
\end{document}
